% Options for packages loaded elsewhere
\PassOptionsToPackage{unicode}{hyperref}
\PassOptionsToPackage{hyphens}{url}
%
\documentclass[
  11pt,
]{article}
\usepackage{amsmath,amssymb}
\usepackage{setspace}
\usepackage{iftex}
\ifPDFTeX
  \usepackage[T1]{fontenc}
  \usepackage[utf8]{inputenc}
  \usepackage{textcomp} % provide euro and other symbols
\else % if luatex or xetex
  \usepackage{unicode-math} % this also loads fontspec
  \defaultfontfeatures{Scale=MatchLowercase}
  \defaultfontfeatures[\rmfamily]{Ligatures=TeX,Scale=1}
\fi
\usepackage{lmodern}
\ifPDFTeX\else
  % xetex/luatex font selection
\fi
% Use upquote if available, for straight quotes in verbatim environments
\IfFileExists{upquote.sty}{\usepackage{upquote}}{}
\IfFileExists{microtype.sty}{% use microtype if available
  \usepackage[]{microtype}
  \UseMicrotypeSet[protrusion]{basicmath} % disable protrusion for tt fonts
}{}
\makeatletter
\@ifundefined{KOMAClassName}{% if non-KOMA class
  \IfFileExists{parskip.sty}{%
    \usepackage{parskip}
  }{% else
    \setlength{\parindent}{0pt}
    \setlength{\parskip}{6pt plus 2pt minus 1pt}}
}{% if KOMA class
  \KOMAoptions{parskip=half}}
\makeatother
\usepackage{xcolor}
\usepackage[top=2.5cm,bottom=1.5cm,left=2.5cm,right=2.0cm]{geometry}
\usepackage{color}
\usepackage{fancyvrb}
\newcommand{\VerbBar}{|}
\newcommand{\VERB}{\Verb[commandchars=\\\{\}]}
\DefineVerbatimEnvironment{Highlighting}{Verbatim}{commandchars=\\\{\}}
% Add ',fontsize=\small' for more characters per line
\usepackage{framed}
\definecolor{shadecolor}{RGB}{248,248,248}
\newenvironment{Shaded}{\begin{snugshade}}{\end{snugshade}}
\newcommand{\AlertTok}[1]{\textcolor[rgb]{0.94,0.16,0.16}{#1}}
\newcommand{\AnnotationTok}[1]{\textcolor[rgb]{0.56,0.35,0.01}{\textbf{\textit{#1}}}}
\newcommand{\AttributeTok}[1]{\textcolor[rgb]{0.13,0.29,0.53}{#1}}
\newcommand{\BaseNTok}[1]{\textcolor[rgb]{0.00,0.00,0.81}{#1}}
\newcommand{\BuiltInTok}[1]{#1}
\newcommand{\CharTok}[1]{\textcolor[rgb]{0.31,0.60,0.02}{#1}}
\newcommand{\CommentTok}[1]{\textcolor[rgb]{0.56,0.35,0.01}{\textit{#1}}}
\newcommand{\CommentVarTok}[1]{\textcolor[rgb]{0.56,0.35,0.01}{\textbf{\textit{#1}}}}
\newcommand{\ConstantTok}[1]{\textcolor[rgb]{0.56,0.35,0.01}{#1}}
\newcommand{\ControlFlowTok}[1]{\textcolor[rgb]{0.13,0.29,0.53}{\textbf{#1}}}
\newcommand{\DataTypeTok}[1]{\textcolor[rgb]{0.13,0.29,0.53}{#1}}
\newcommand{\DecValTok}[1]{\textcolor[rgb]{0.00,0.00,0.81}{#1}}
\newcommand{\DocumentationTok}[1]{\textcolor[rgb]{0.56,0.35,0.01}{\textbf{\textit{#1}}}}
\newcommand{\ErrorTok}[1]{\textcolor[rgb]{0.64,0.00,0.00}{\textbf{#1}}}
\newcommand{\ExtensionTok}[1]{#1}
\newcommand{\FloatTok}[1]{\textcolor[rgb]{0.00,0.00,0.81}{#1}}
\newcommand{\FunctionTok}[1]{\textcolor[rgb]{0.13,0.29,0.53}{\textbf{#1}}}
\newcommand{\ImportTok}[1]{#1}
\newcommand{\InformationTok}[1]{\textcolor[rgb]{0.56,0.35,0.01}{\textbf{\textit{#1}}}}
\newcommand{\KeywordTok}[1]{\textcolor[rgb]{0.13,0.29,0.53}{\textbf{#1}}}
\newcommand{\NormalTok}[1]{#1}
\newcommand{\OperatorTok}[1]{\textcolor[rgb]{0.81,0.36,0.00}{\textbf{#1}}}
\newcommand{\OtherTok}[1]{\textcolor[rgb]{0.56,0.35,0.01}{#1}}
\newcommand{\PreprocessorTok}[1]{\textcolor[rgb]{0.56,0.35,0.01}{\textit{#1}}}
\newcommand{\RegionMarkerTok}[1]{#1}
\newcommand{\SpecialCharTok}[1]{\textcolor[rgb]{0.81,0.36,0.00}{\textbf{#1}}}
\newcommand{\SpecialStringTok}[1]{\textcolor[rgb]{0.31,0.60,0.02}{#1}}
\newcommand{\StringTok}[1]{\textcolor[rgb]{0.31,0.60,0.02}{#1}}
\newcommand{\VariableTok}[1]{\textcolor[rgb]{0.00,0.00,0.00}{#1}}
\newcommand{\VerbatimStringTok}[1]{\textcolor[rgb]{0.31,0.60,0.02}{#1}}
\newcommand{\WarningTok}[1]{\textcolor[rgb]{0.56,0.35,0.01}{\textbf{\textit{#1}}}}
\usepackage{longtable,booktabs,array}
\usepackage{calc} % for calculating minipage widths
% Correct order of tables after \paragraph or \subparagraph
\usepackage{etoolbox}
\makeatletter
\patchcmd\longtable{\par}{\if@noskipsec\mbox{}\fi\par}{}{}
\makeatother
% Allow footnotes in longtable head/foot
\IfFileExists{footnotehyper.sty}{\usepackage{footnotehyper}}{\usepackage{footnote}}
\makesavenoteenv{longtable}
\setlength{\emergencystretch}{3em} % prevent overfull lines
\providecommand{\tightlist}{%
  \setlength{\itemsep}{0pt}\setlength{\parskip}{0pt}}
\setcounter{secnumdepth}{5}
% definitions for citeproc citations
\NewDocumentCommand\citeproctext{}{}
\NewDocumentCommand\citeproc{mm}{%
  \begingroup\def\citeproctext{#2}\cite{#1}\endgroup}
\makeatletter
 % allow citations to break across lines
 \let\@cite@ofmt\@firstofone
 % avoid brackets around text for \cite:
 \def\@biblabel#1{}
 \def\@cite#1#2{{#1\if@tempswa , #2\fi}}
\makeatother
\newlength{\cslhangindent}
\setlength{\cslhangindent}{1.5em}
\newlength{\csllabelwidth}
\setlength{\csllabelwidth}{3em}
\newenvironment{CSLReferences}[2] % #1 hanging-indent, #2 entry-spacing
 {\begin{list}{}{%
  \setlength{\itemindent}{0pt}
  \setlength{\leftmargin}{0pt}
  \setlength{\parsep}{0pt}
  % turn on hanging indent if param 1 is 1
  \ifodd #1
   \setlength{\leftmargin}{\cslhangindent}
   \setlength{\itemindent}{-1\cslhangindent}
  \fi
  % set entry spacing
  \setlength{\itemsep}{#2\baselineskip}}}
 {\end{list}}
\usepackage{calc}
\newcommand{\CSLBlock}[1]{\hfill\break\parbox[t]{\linewidth}{\strut\ignorespaces#1\strut}}
\newcommand{\CSLLeftMargin}[1]{\parbox[t]{\csllabelwidth}{\strut#1\strut}}
\newcommand{\CSLRightInline}[1]{\parbox[t]{\linewidth - \csllabelwidth}{\strut#1\strut}}
\newcommand{\CSLIndent}[1]{\hspace{\cslhangindent}#1}
\usepackage{booktabs}
\usepackage{fontspec}
\usepackage{amsmath}
\usepackage{amssymb}
\usepackage{tabularx}
\usepackage{fancyhdr}
\IfFontExistsTF{Arial}{\newfontfamily\headerfont{Arial}}{\newfontfamily\headerfont{Arimo}}
\setlength{\headheight}{14pt}
\AtBeginDocument{
  \pagestyle{fancy}
  \fancyhf{}
  \fancyhead[L]{\headerfont\fontsize{9}{11}\selectfont Segment Platform Research Foundation Document}
  \fancyhead[R]{\headerfont\fontsize{9}{11}\selectfont \thepage}
  \renewcommand{\headrulewidth}{0pt}
  \renewcommand{\footrulewidth}{0pt}
  \makeatletter\let\ps@plain\ps@fancy\makeatother
}
\makeatletter
\@ifpackageloaded{subfig}{}{\usepackage{subfig}}
\@ifpackageloaded{caption}{}{\usepackage{caption}}
\captionsetup[subfloat]{margin=0.5em}
\AtBeginDocument{%
\renewcommand*\figurename{Figure}
\renewcommand*\tablename{Table}
}
\AtBeginDocument{%
\renewcommand*\listfigurename{List of Figures}
\renewcommand*\listtablename{List of Tables}
}
\newcounter{pandoccrossref@subfigures@footnote@counter}
\newenvironment{pandoccrossrefsubfigures}{%
\setcounter{pandoccrossref@subfigures@footnote@counter}{0}
\begin{figure}\centering%
\gdef\global@pandoccrossref@subfigures@footnotes{}%
\DeclareRobustCommand{\footnote}[1]{\footnotemark%
\stepcounter{pandoccrossref@subfigures@footnote@counter}%
\ifx\global@pandoccrossref@subfigures@footnotes\empty%
\gdef\global@pandoccrossref@subfigures@footnotes{{##1}}%
\else%
\g@addto@macro\global@pandoccrossref@subfigures@footnotes{, {##1}}%
\fi}}%
{\end{figure}%
\addtocounter{footnote}{-\value{pandoccrossref@subfigures@footnote@counter}}
\@for\f:=\global@pandoccrossref@subfigures@footnotes\do{\stepcounter{footnote}\footnotetext{\f}}%
\gdef\global@pandoccrossref@subfigures@footnotes{}}
\@ifpackageloaded{float}{}{\usepackage{float}}
\floatstyle{ruled}
\@ifundefined{c@chapter}{\newfloat{codelisting}{h}{lop}}{\newfloat{codelisting}{h}{lop}[chapter]}
\floatname{codelisting}{Listing}
\newcommand*\listoflistings{\listof{codelisting}{List of Listings}}
\makeatother
\ifLuaTeX
  \usepackage{selnolig}  % disable illegal ligatures
\fi
\usepackage{bookmark}
\IfFileExists{xurl.sty}{\usepackage{xurl}}{} % add URL line breaks if available
\urlstyle{same}
\hypersetup{
  pdftitle={Segment Platform Research Foundation Document},
  pdfauthor={Steffen Lukas steffen.lukas@charite.de},
  hidelinks,
  pdfcreator={LaTeX via pandoc}}

\title{Segment Platform Research Foundation Document}
\usepackage{etoolbox}
\makeatletter
\providecommand{\subtitle}[1]{% add subtitle to \maketitle
  \apptocmd{\@title}{\par {\large #1 \par}}{}{}
}
\makeatother
\subtitle{Browser-Based Interactive 3D Medical Segmentation with
SAM-Med3D-turbo}
\author{Steffen Lukas
\href{mailto:steffen.lukas@charite.de}{\nolinkurl{steffen.lukas@charite.de}}}
\date{13 March 2026}

\begin{document}
\maketitle

{
\setcounter{tocdepth}{2}
\tableofcontents
}
\setstretch{1.2}
\subsection*{Abstract}\label{abstract}
\addcontentsline{toc}{subsection}{Abstract}

This document is a technical and research foundation for the Segment
platform, a browser-based 3D medical segmentation environment built
around SAM-Med3D-turbo. It defines the clinical targets, model
assumptions, workflow patterns, system architecture, and evaluation path
for two initial application domains: coronary CCTA in DISCHARGE and
prostate mpMRI. It is an internal technical reference and research
planning document rather than a deployment certification package.

\subsection*{Purpose}\label{purpose}
\addcontentsline{toc}{subsection}{Purpose}

This reference serves as the working document for framing the
segmentation problem, the platform requirements, and the validation
roadmap for Charité-led development.

It is \textbf{not}:

\begin{itemize}
\tightlist
\item
  a production deployment architecture specification;
\item
  a finished browser-platform product brief;
\item
  a standalone funding proposal.
\end{itemize}

\subsection*{Why Segmentation Matters
Here}\label{why-segmentation-matters-here}
\addcontentsline{toc}{subsection}{Why Segmentation Matters Here}

The current programme needs segmentation for three distinct but linked
tasks:

\begin{enumerate}
\def\labelenumi{\arabic{enumi}.}
\tightlist
\item
  \textbf{Coronary lumen geometry} for stenosis quantification,
  centreline extraction, and straightened MPR analysis.
\item
  \textbf{Outer-wall and plaque context} where vessel-wall estimation
  and plaque-burden measurements are clinically relevant.
\item
  \textbf{Prostate zones, lesions, and organs at risk} for biopsy
  support, radiotherapy planning, and multi-domain testing of
  foundation-model generalisation.
\end{enumerate}

\subsection*{Document Snapshot}\label{document-snapshot}
\addcontentsline{toc}{subsection}{Document Snapshot}

\begin{longtable}[]{@{}
  >{\raggedright\arraybackslash}p{(\columnwidth - 2\tabcolsep) * \real{0.4667}}
  >{\raggedright\arraybackslash}p{(\columnwidth - 2\tabcolsep) * \real{0.5333}}@{}}
\toprule\noalign{}
\begin{minipage}[b]{\linewidth}\raggedright
Field
\end{minipage} & \begin{minipage}[b]{\linewidth}\raggedright
Detail
\end{minipage} \\
\midrule\noalign{}
\endhead
\bottomrule\noalign{}
\endlastfoot
\textbf{Institution} & Charité - Universitätsmedizin Berlin \\
\textbf{Core Stack} & TypeScript · Vite · Niivue 0.66 · FastAPI ·
PyTorch \\
\textbf{Foundation Model} & SAM-Med3D-turbo (3D ViT-B/16, 91 M
params) \\
\textbf{Primary Dataset} & DISCHARGE (\textasciitilde25 M DICOM slices
across 3,561 patients) \\
\textbf{Pre-training Corpus} & SA-Med3D-140K (21 729 volumes, 143 518
masks, 245 categories) \\
\textbf{Clinical Domains} & Coronary CCTA (lumen / outer wall / plaque)
· Prostate mpMRI (zones / lesions / OARs) \\
\end{longtable}

\subsection{Executive Summary \& Research
Vision}\label{executive-summary-research-vision}

\subsubsection{Clinical Gap}\label{clinical-gap}

\begin{longtable}[]{@{}
  >{\raggedright\arraybackslash}p{(\columnwidth - 2\tabcolsep) * \real{0.5294}}
  >{\raggedright\arraybackslash}p{(\columnwidth - 2\tabcolsep) * \real{0.4706}}@{}}
\toprule\noalign{}
\begin{minipage}[b]{\linewidth}\raggedright
Problem
\end{minipage} & \begin{minipage}[b]{\linewidth}\raggedright
Impact
\end{minipage} \\
\midrule\noalign{}
\endhead
\bottomrule\noalign{}
\endlastfoot
Manual coronary segmentation & 30--60 min / case, €200--400 \\
Limited scalability & DISCHARGE (3 561 patients) still largely
un-segmented \\
Prostate mpMRI zone/lesion contouring & Equally time-intensive for
biopsy/radiotherapy planning \\
Centre-specific expertise & Manual segmentation available only at
specialised sites \\
\end{longtable}

\subsubsection{Our Solution}\label{our-solution}

A \textbf{single, universal, browser-based platform} powered by
\textbf{SAM-Med3D-turbo}:

\begin{itemize}
\tightlist
\item
  \textbf{1--3 point prompts} (3D coordinates in mm space) → any
  anatomy, any modality
\item
  \textbf{Zero installation} --- runs in Chrome / Firefox / Safari
\item
  \textbf{On-premise GPU inference} at Charité (GDPR-compliant; no data
  leaves the hospital)
\item
  \textbf{Interactive workflow} for radiologists + \textbf{batch
  research mode} for DISCHARGE processing + \textbf{active-learning
  loop} for continuous improvement
\end{itemize}

\subsubsection{Research Foundation}\label{research-foundation}

This platform is a \textbf{test-bed for foundation-model research} in
cardiovascular imaging:

\begin{enumerate}
\def\labelenumi{\arabic{enumi}.}
\tightlist
\item
  Evaluate \textbf{zero-shot / few-shot generalisation} of
  SAM-Med3D-turbo on unseen DISCHARGE cases
\item
  Quantify \textbf{active-learning gains} (weekly fine-tuning on expert
  corrections)
\item
  Benchmark against \textbf{nnU-Net} task-specific models on clinical
  endpoints (stenosis \%, plaque burden, FFR-CT correlation)
\item
  Open-source modular components for the broader MedAI community
\end{enumerate}

\subsubsection{Success Targets (6-month
horizon)}\label{success-targets-6-month-horizon}

\begin{longtable}[]{@{}ll@{}}
\toprule\noalign{}
Metric & Target \\
\midrule\noalign{}
\endhead
\bottomrule\noalign{}
\endlastfoot
DISCHARGE cases auto-processed & \textgreater{} 80 \% of dataset \\
End-to-end latency & \textless{} 2 s per vessel / per prostate gland \\
Dice (coronary lumen) & \textgreater{} 0.85 vs.~expert \\
Dice (prostate whole-gland) & \textgreater{} 0.90 vs.~expert \\
Cost reduction & 10× cheaper than manual contouring \\
\end{longtable}

\begin{quote}
\textbf{Critical note on Dice targets:} The SAM-Med3D paper reports
\textbf{87.12 \% Dice on cardiac structures} with 1 prompt point (Table
5 in\textsuperscript{1}). However, this was measured on the ACDC dataset
(cardiac MRI short-axis cine), \textbf{not} coronary CTA. Coronary
arteries are smaller, noisier, and motion-affected --- published
coronary lumen Dice values for task-specific models
(nnU-Net\textsuperscript{2}) range 0.75--0.88 depending on vessel
branch. Our 0.85 target is therefore ambitious but grounded.
\end{quote}

\begin{quote}
\textbf{Hardware \& deployment requirements} (GPU specs, production
cluster, GDPR constraints) are in §13.
\end{quote}

\begin{center}\rule{0.5\linewidth}{0.5pt}\end{center}

\subsection{Foundation Model: SAM-Med3D-turbo --- Verified Technical
Profile}\label{foundation-model-sam-med3d-turbo-verified-technical-profile}

\subsubsection{Paper \& Publication}\label{paper-publication}

\begin{longtable}[]{@{}
  >{\raggedright\arraybackslash}p{(\columnwidth - 2\tabcolsep) * \real{0.5000}}
  >{\raggedright\arraybackslash}p{(\columnwidth - 2\tabcolsep) * \real{0.5000}}@{}}
\toprule\noalign{}
\begin{minipage}[b]{\linewidth}\raggedright
Field
\end{minipage} & \begin{minipage}[b]{\linewidth}\raggedright
Value
\end{minipage} \\
\midrule\noalign{}
\endhead
\bottomrule\noalign{}
\endlastfoot
\textbf{Title} & SAM-Med3D: Towards General-purpose Segmentation Models
for Volumetric Medical Images \\
\textbf{Authors} & Haoyu Wang, Sizheng Guo, Jin Ye, Zhongying Deng,
Junlong Cheng, Tianbin Li, Jianpin Chen, Yanzhou Su, Ziyan Huang, Yiqing
Shen, Bin Fu, Shaoting Zhang, Junjun He, Yu Qiao \\
\textbf{Venue} & \textbf{ECCV BIC 2024 --- Oral} \\
\textbf{arXiv} & \href{https://arxiv.org/abs/2310.15161}{2310.15161} \\
\textbf{License} & Apache 2.0 \\
\end{longtable}

\subsubsection{Architecture (Verified from Paper §4.1 \&
GitHub)}\label{architecture-verified-from-paper-4.1-github}

SAM-Med3D uses a \textbf{fully 3D architecture trained from scratch}
(Method 3 in the paper). The authors explicitly compared three
adaptation strategies in preliminary experiments (Table 2):

\begin{longtable}[]{@{}llll@{}}
\toprule\noalign{}
Strategy & Seen Dice & Unseen Dice & Chosen? \\
\midrule\noalign{}
\endhead
\bottomrule\noalign{}
\endlastfoot
3D Adapter + Frozen SAM & Lower & Moderate & No \\
Fine-tune SAM 2D→3D weights & Good & Poor (broken priors) & No \\
\textbf{Train 3D from scratch} & \textbf{Good} & \textbf{Best} & Yes \\
\end{longtable}

\textbf{Rationale (Paper §4.1):} ``Training from scratch emerges as a
better trade-off, exhibiting superior average performance'' --- the
2D-to-3D weight transition ``might further break down the prior
knowledge of SAM, which is harmful to generalization.''

\textbf{Component breakdown:}

\begin{longtable}[]{@{}
  >{\raggedright\arraybackslash}p{(\columnwidth - 4\tabcolsep) * \real{0.3056}}
  >{\raggedright\arraybackslash}p{(\columnwidth - 4\tabcolsep) * \real{0.3611}}
  >{\raggedright\arraybackslash}p{(\columnwidth - 4\tabcolsep) * \real{0.3333}}@{}}
\toprule\noalign{}
\begin{minipage}[b]{\linewidth}\raggedright
Component
\end{minipage} & \begin{minipage}[b]{\linewidth}\raggedright
Architecture
\end{minipage} & \begin{minipage}[b]{\linewidth}\raggedright
Parameters
\end{minipage} \\
\midrule\noalign{}
\endhead
\bottomrule\noalign{}
\endlastfoot
Image Encoder & 3D ViT-B/16 (3D positional encoding, 3D convolutions, 3D
LayerNorm) & \textasciitilde86 M \\
Prompt Encoder & 3D point/box encoder (learned embeddings) &
\textasciitilde1 M \\
Mask Decoder & Lightweight 3D decoder (2-layer transformer + upsampling)
& \textasciitilde4 M \\
\textbf{Total} & & \textbf{\textasciitilde91 M} \\
\end{longtable}

\begin{quote}
\textbf{Critical note:} The paper states ``86M encoder + 5M decoder'' in
various summaries. The exact split varies by source. The model is
\textbf{not} a modified SAM (Meta) --- it is architecturally distinct,
sharing only the conceptual prompt-based paradigm.
\end{quote}

\subsubsection{Training (Verified from Paper
§4.2)}\label{training-verified-from-paper-4.2}

\textbf{Two-stage procedure:}

\begin{longtable}[]{@{}
  >{\raggedright\arraybackslash}p{(\columnwidth - 6\tabcolsep) * \real{0.2333}}
  >{\raggedright\arraybackslash}p{(\columnwidth - 6\tabcolsep) * \real{0.2000}}
  >{\raggedright\arraybackslash}p{(\columnwidth - 6\tabcolsep) * \real{0.2667}}
  >{\raggedright\arraybackslash}p{(\columnwidth - 6\tabcolsep) * \real{0.3000}}@{}}
\toprule\noalign{}
\begin{minipage}[b]{\linewidth}\raggedright
Stage
\end{minipage} & \begin{minipage}[b]{\linewidth}\raggedright
Data
\end{minipage} & \begin{minipage}[b]{\linewidth}\raggedright
Epochs
\end{minipage} & \begin{minipage}[b]{\linewidth}\raggedright
Purpose
\end{minipage} \\
\midrule\noalign{}
\endhead
\bottomrule\noalign{}
\endlastfoot
\textbf{Stage 1: Pre-training} & All 131 K masks from SA-Med3D-140K
training set & 800 & Build general 3D medical understanding \\
\textbf{Stage 2: Fine-tuning} & \textasciitilde75 K filtered
high-quality masks & Additional & Improve prompt efficiency on
challenging targets \\
\end{longtable}

\textbf{SAM-Med3D-turbo} (from
\href{https://github.com/uni-medical/SAM-Med3D/issues/2\#issuecomment-1849002225}{GitHub
issue \#2}): - Fine-tuned on \textbf{44 additional datasets} beyond the
base SA-Med3D-140K - Optimised for \textbf{sub-second inference} with
FP16 - This is the recommended checkpoint for deployment

\subsubsection{SA-Med3D-140K Dataset (Verified from
HuggingFace)}\label{sa-med3d-140k-dataset-verified-from-huggingface}

\begin{longtable}[]{@{}
  >{\raggedright\arraybackslash}p{(\columnwidth - 2\tabcolsep) * \real{0.6111}}
  >{\raggedright\arraybackslash}p{(\columnwidth - 2\tabcolsep) * \real{0.3889}}@{}}
\toprule\noalign{}
\begin{minipage}[b]{\linewidth}\raggedright
Statistic
\end{minipage} & \begin{minipage}[b]{\linewidth}\raggedright
Value
\end{minipage} \\
\midrule\noalign{}
\endhead
\bottomrule\noalign{}
\endlastfoot
Total 3D images & 21 729 \\
Total 3D masks & 143 518 \\
Anatomical categories & 245 \\
Modalities & 28 (CT, MR, US, and more) \\
Sources & 70 public datasets + 8,128 privately licensed cases from 24
hospitals \\
Primary task & General-purpose promptable segmentation \\
\end{longtable}

\subsubsection{Verified Performance (from Paper Tables
3--5)}\label{verified-performance-from-paper-tables-35}

\textbf{Overall (Table 3):} - SAM-Med3D with 1 point: \textbf{+60.12 \%
overall Dice} improvement over original SAM - Consistently outperforms
SAM-Med2D across all prompt counts - Operates at \textbf{1--26 \% of
inference time} compared to slice-by-slice SAM

\textbf{By anatomy (Table 5, 1 prompt point):}

\begin{longtable}[]{@{}llll@{}}
\toprule\noalign{}
Anatomy & SAM & SAM-Med2D & SAM-Med3D \\
\midrule\noalign{}
\endhead
\bottomrule\noalign{}
\endlastfoot
Cardiac (seen) & Poor & Moderate & \textbf{87.12 \%} \\
Organs (seen) & Poor & Moderate & \textbf{Best} \\
Lesions (unseen) & Very poor & Moderate & Competitive \\
Bones/muscles & Very poor & Moderate & \textbf{Greatest advantage} \\
\end{longtable}

\textbf{Key finding (Paper §5.1.5):} ``SAM-Med3D using 1 point
outperforms SAM-Med2D with N points in \textbf{45 targets out of 49},
achieving up to +68.2\% improvement.''

\textbf{Transferability (Paper §5.1.6, Table 6):} When SAM-Med3D's ViT
encoder is used as pre-trained backbone for UNETR, it yields up to
\textbf{+5.63 \% Dice improvement} over training from scratch ---
confirming value as a foundation model.

\begin{quote}
\textbf{Critical caveat for our project:} The ``cardiac'' results (87.12
\%) are from the \textbf{ACDC dataset} (cardiac MRI, short-axis cine ---
segmenting LV/RV/myocardium). This is \textbf{not} coronary artery
segmentation from CTA. Coronary arteries are 1.5--4 mm diameter,
motion-affected, and require dual-wall segmentation --- a substantially
harder task not directly evaluated in the paper. The model has
\textbf{never been benchmarked on coronary CTA lumen segmentation}. Our
project will provide this missing evaluation.
\end{quote}

\begin{quote}
\textbf{Strategic implication:} Because LV/RV/myocardium from cardiac
MRI \emph{is} the dominant pre-trained cardiac target,
\textbf{myocardial segmentation from CCTA should be the first zero-shot
validation task} --- not coronary lumen. It is the easiest
coronary-domain target for the model given pre-training priors, provides
a large structure (easy Dice baseline), and is directly needed for
FFR-CT territory assignment and MACE prediction from the DISCHARGE data.
The coronary lumen evaluation should follow once the myocardial baseline
is established.
\end{quote}

\subsubsection{Official Resources}\label{official-resources}

\begin{longtable}[]{@{}
  >{\raggedright\arraybackslash}p{(\columnwidth - 2\tabcolsep) * \real{0.6250}}
  >{\raggedright\arraybackslash}p{(\columnwidth - 2\tabcolsep) * \real{0.3750}}@{}}
\toprule\noalign{}
\begin{minipage}[b]{\linewidth}\raggedright
Resource
\end{minipage} & \begin{minipage}[b]{\linewidth}\raggedright
Link
\end{minipage} \\
\midrule\noalign{}
\endhead
\bottomrule\noalign{}
\endlastfoot
GitHub &
\href{https://github.com/uni-medical/SAM-Med3D}{uni-medical/SAM-Med3D} \\
Paper & \href{https://arxiv.org/abs/2310.15161}{arXiv:2310.15161} \\
Supplementary &
\href{https://github.com/uni-medical/SAM-Med3D/blob/main/paper/SAM_Med3D_ECCV_Supplementary.pdf}{ECCV
Supplementary PDF} \\
Turbo checkpoint &
\href{https://huggingface.co/blueyo0/SAM-Med3D/resolve/main/sam_med3d_turbo.pth}{HuggingFace:
sam\_med3d\_turbo.pth} \\
Dataset &
\href{https://huggingface.co/datasets/blueyo0/SA-Med3D-140K}{HuggingFace:
SA-Med3D-140K} \\
MedIM loader &
\href{https://github.com/uni-medical/MedIM}{uni-medical/MedIM} \\
CVPR25 Challenge &
\href{https://www.codabench.org/competitions/5263/}{MedSegFM
Competition} \\
\end{longtable}

\subsubsection{Environment Setup (Verified from GitHub
README)}\label{environment-setup-verified-from-github-readme}

\textbf{Model / research environment:}

\begin{Shaded}
\begin{Highlighting}[]
\ExtensionTok{conda}\NormalTok{ create }\AttributeTok{{-}{-}name}\NormalTok{ sammed3d python=3.10}
\ExtensionTok{conda}\NormalTok{ activate sammed3d}
\ExtensionTok{pip}\NormalTok{ install uv}
\ExtensionTok{uv}\NormalTok{ pip install torch==2.6.0 torchvision==0.21.0 torchaudio==2.6.0}
\ExtensionTok{uv}\NormalTok{ pip install torchio opencv{-}python{-}headless matplotlib prefetch\_generator monai edt surface{-}distance medim}
\end{Highlighting}
\end{Shaded}

\textbf{Backend (FastAPI server) --- additional dependencies:}

\begin{Shaded}
\begin{Highlighting}[]
\ExtensionTok{uv}\NormalTok{ pip install fastapi uvicorn nibabel SimpleITK celery redis}
\end{Highlighting}
\end{Shaded}

\textbf{Compile this document to PDF:}

\begin{Shaded}
\begin{Highlighting}[]
\ExtensionTok{./code/md2pdf.sh}\NormalTok{ plan/sam3d.md}
\CommentTok{\# Optional: also produce DOCX}
\ExtensionTok{./code/md2pdf.sh}\NormalTok{ plan/sam3d.md }\AttributeTok{{-}{-}docx}
\end{Highlighting}
\end{Shaded}

\subsubsection{Model Loading (Verified from
GitHub)}\label{model-loading-verified-from-github}

\begin{Shaded}
\begin{Highlighting}[]
\ImportTok{import}\NormalTok{ medim}

\CommentTok{\# Option A: Direct from HuggingFace (downloads automatically)}
\CommentTok{\# Use /resolve/ (raw file), NOT /blob/ (HTML page)}
\NormalTok{ckpt\_path }\OperatorTok{=} \StringTok{"https://huggingface.co/blueyo0/SAM{-}Med3D/resolve/main/sam\_med3d\_turbo.pth"}
\NormalTok{model }\OperatorTok{=}\NormalTok{ medim.create\_model(}\StringTok{"SAM{-}Med3D"}\NormalTok{, pretrained}\OperatorTok{=}\VariableTok{True}\NormalTok{, checkpoint\_path}\OperatorTok{=}\NormalTok{ckpt\_path)}

\CommentTok{\# Option B: Local checkpoint (recommended for deployment)}
\NormalTok{model }\OperatorTok{=}\NormalTok{ medim.create\_model(}
    \StringTok{"SAM{-}Med3D"}\NormalTok{,}
\NormalTok{    pretrained}\OperatorTok{=}\VariableTok{True}\NormalTok{,}
\NormalTok{    checkpoint\_path}\OperatorTok{=}\StringTok{"app/models/checkpoints/sam\_med3d\_turbo.pth"}
\NormalTok{)}

\CommentTok{\# Optimise for inference}
\ImportTok{import}\NormalTok{ torch}
\NormalTok{device }\OperatorTok{=} \StringTok{"cuda"} \ControlFlowTok{if}\NormalTok{ torch.cuda.is\_available() }\ControlFlowTok{else} \StringTok{"cpu"}
\NormalTok{model }\OperatorTok{=}\NormalTok{ model.to(device).half()  }\CommentTok{\# FP16 for speed + memory}
\NormalTok{model.}\BuiltInTok{eval}\NormalTok{()}
\end{Highlighting}
\end{Shaded}

\subsubsection{Data Format for Fine-tuning (Verified from
GitHub)}\label{data-format-for-fine-tuning-verified-from-github}

SAM-Med3D expects \textbf{nnU-Net-style}\textsuperscript{2} directory
layout with \textbf{per-structure binary masks}. Unlike nnU-Net which
accepts a single multi-class integer mask (0=background, 1=class1,
\ldots), SAM-Med3D requires one separate binary NIfTI file (values 0/1)
per anatomical structure. Using a single multi-class mask will silently
produce incorrect training prompts.

\begin{verbatim}
data/medical_preprocessed/
├── coronary/
│   ├── ct_DISCHARGE/
│   │   ├── imagesTr/
│   │   │   ├── discharge_0001.nii.gz
│   │   │   └── ...
│   │   └── labelsTr/
│   │       ├── discharge_0001.nii.gz  (binary mask)
│   │       └── ...
└── prostate/
    └── mr_CHARITE/
        ├── imagesTr/
        └── labelsTr/
\end{verbatim}

\begin{quote}
\textbf{Important (from GitHub):} Ground-truth labels are required to
generate prompt points during training. For inference without ground
truth, ``generate a fake ground-truth with the target region for prompt
annotated.''
\end{quote}

\subsubsection{Turbo vs.~Standard
Comparison}\label{turbo-vs.-standard-comparison}

\begin{longtable}[]{@{}lll@{}}
\toprule\noalign{}
Metric & Standard (base .pth) & Turbo (sam\_med3d\_turbo.pth) \\
\midrule\noalign{}
\endhead
\bottomrule\noalign{}
\endlastfoot
Pre-training data & 131 K masks & 131 K + 44 additional datasets \\
Average Dice (1 prompt) & \textasciitilde0.75 & \textasciitilde0.82+ \\
Inference time & 4--8 s & 0.5--1.5 s \\
VRAM usage & \textasciitilde10 GB (FP32) & \textasciitilde4 GB (FP16) \\
\end{longtable}

\begin{center}\rule{0.5\linewidth}{0.5pt}\end{center}

\subsection{Clinical Domain A: Coronary CCTA Segmentation
(DISCHARGE)}\label{clinical-domain-a-coronary-ccta-segmentation-discharge}

\subsubsection{DISCHARGE Trial Overview}\label{discharge-trial-overview}

\begin{longtable}[]{@{}
  >{\raggedright\arraybackslash}p{(\columnwidth - 2\tabcolsep) * \real{0.5000}}
  >{\raggedright\arraybackslash}p{(\columnwidth - 2\tabcolsep) * \real{0.5000}}@{}}
\toprule\noalign{}
\begin{minipage}[b]{\linewidth}\raggedright
Field
\end{minipage} & \begin{minipage}[b]{\linewidth}\raggedright
Value
\end{minipage} \\
\midrule\noalign{}
\endhead
\bottomrule\noalign{}
\endlastfoot
\textbf{Full name} & Diagnostic Imaging Strategies for Patients with
Stable Chest Pain and Intermediate Risk of Coronary Artery Disease \\
\textbf{Reference} & Dewey et al., NEJM 2022\textsuperscript{3} \\
\textbf{Design} & Multicentre randomised controlled trial \\
\textbf{Patients} & 3 561 \\
\textbf{Images} & \textasciitilde25 M DICOM slices (across 3,561
patients; not 25 M cases) \\
\textbf{Modality} & Coronary Computed Tomography Angiography (CCTA) \\
\textbf{Institution} & Charité, Berlin (lead site) \\
\end{longtable}

\subsubsection{\texorpdfstring{Standardised Nomenclature (CAD-RADS 2.0 /
AHA 17-segment
compatible)\textsuperscript{4}}{Standardised Nomenclature (CAD-RADS 2.0 / AHA 17-segment compatible)4}}\label{standardised-nomenclature-cad-rads-2.0-aha-17-segment-compatible-maurovich2020cadrads}

\begin{longtable}[]{@{}
  >{\raggedright\arraybackslash}p{(\columnwidth - 6\tabcolsep) * \real{0.1954}}
  >{\raggedright\arraybackslash}p{(\columnwidth - 6\tabcolsep) * \real{0.2989}}
  >{\raggedright\arraybackslash}p{(\columnwidth - 6\tabcolsep) * \real{0.1494}}
  >{\raggedright\arraybackslash}p{(\columnwidth - 6\tabcolsep) * \real{0.3563}}@{}}
\toprule\noalign{}
\begin{minipage}[b]{\linewidth}\raggedright
Clinical Feature
\end{minipage} & \begin{minipage}[b]{\linewidth}\raggedright
Segmentation Class Label
\end{minipage} & \begin{minipage}[b]{\linewidth}\raggedright
Description
\end{minipage} & \begin{minipage}[b]{\linewidth}\raggedright
HU Range (contrast-enhanced CT)
\end{minipage} \\
\midrule\noalign{}
\endhead
\bottomrule\noalign{}
\endlastfoot
\textbf{Myocardium (LV)} & \texttt{LV\_MYO} & Muscular wall of left
ventricle & 50--120 HU \\
\textbf{Endocardial Lumen} & \texttt{Endocardial\_Lumen} & Inner chamber
blood pool & 250--400 HU \\
\textbf{Coronary Lumen} &
\texttt{Lumen\_\{LAD\textbackslash{}\textbar{}RCA\textbackslash{}\textbar{}LCx\textbackslash{}\textbar{}LM\}}
& Vessel blood pool --- per AHA branch & 300--400 HU \\
\textbf{Outer Wall (EEM)} & \texttt{VesselWall\_EEM} & External elastic
membrane boundary & --- (morphological) \\
\textbf{Calcified Plaque} & \texttt{Plaque\_Calc} & High-density stable
plaque & \textgreater{} 130 HU (often 500--1000) \\
\textbf{Fibrous Plaque} & \texttt{Plaque\_Fibrous} &
Intermediate-density plaque & 60--130 HU \\
\textbf{Low-Attenuation Plaque} & \texttt{Plaque\_LAP} & Lipid-rich /
necrotic core --- \textbf{high risk} & \textless{} 60 HU \\
\end{longtable}

\subsubsection{Why Coronary Arteries are the Hardest Segmentation
Target}\label{why-coronary-arteries-are-the-hardest-segmentation-target}

\begin{longtable}[]{@{}
  >{\raggedright\arraybackslash}p{(\columnwidth - 4\tabcolsep) * \real{0.4074}}
  >{\raggedright\arraybackslash}p{(\columnwidth - 4\tabcolsep) * \real{0.2963}}
  >{\raggedright\arraybackslash}p{(\columnwidth - 4\tabcolsep) * \real{0.2963}}@{}}
\toprule\noalign{}
\begin{minipage}[b]{\linewidth}\raggedright
Challenge
\end{minipage} & \begin{minipage}[b]{\linewidth}\raggedright
Detail
\end{minipage} & \begin{minipage}[b]{\linewidth}\raggedright
Impact
\end{minipage} \\
\midrule\noalign{}
\endhead
\bottomrule\noalign{}
\endlastfoot
\textbf{Motion artefacts} & Heart beats 60--80 bpm; RCA most affected &
Blurred vessel edges despite ECG gating \\
\textbf{Small calibre} & 1.5--4 mm diameter & Partial volume effects at
0.5 mm resolution \\
\textbf{Low-contrast plaque} & Lipid-rich plaque 20--50 HU vs.~lumen
300--400 HU & Nearly invisible on standard windowing \\
\textbf{Blooming} & Calcification \textgreater{} 130 HU causes beam
hardening & Obscures adjacent soft plaque \\
\textbf{Complex topology} & Bifurcations, overlapping branches,
tortuosity & Single-click prompts often insufficient \\
\textbf{Dual-wall requirement} & Must segment lumen AND outer wall
separately & Wall thickness 0.5--3 mm = 1--5 voxels \\
\end{longtable}

\subsubsection{SAM-Med3D-turbo Prompting Strategy for Coronary
CCTA}\label{sam-med3d-turbo-prompting-strategy-for-coronary-ccta}

\paragraph{A. Multi-Point Centerline Prompt (``String''
Prompt)}\label{a.-multi-point-centerline-prompt-string-prompt}

Standard single-click prompts fail for thin, tortuous vessels spanning
100+ slices. Instead:

\begin{enumerate}
\def\labelenumi{\arabic{enumi}.}
\tightlist
\item
  User clicks \textbf{proximal ostium} (positive point)
\item
  User clicks \textbf{distal vessel tip} (positive point)
\item
  Model ``fills in'' the lumen tube between points
\item
  If mask leaks into \textbf{great cardiac vein} → place
  \textbf{negative point} on vein
\end{enumerate}

\begin{Shaded}
\begin{Highlighting}[]
\CommentTok{\# Coronary lumen segmentation with multi{-}point prompt}
\NormalTok{lumen\_mask }\OperatorTok{=}\NormalTok{ model.segment(}
\NormalTok{    volume}\OperatorTok{=}\NormalTok{cta\_volume,}
\NormalTok{    points}\OperatorTok{=}\NormalTok{[ostium\_xyz, mid\_vessel\_xyz, distal\_xyz],}
\NormalTok{    labels}\OperatorTok{=}\NormalTok{[}\DecValTok{1}\NormalTok{, }\DecValTok{1}\NormalTok{, }\DecValTok{1}\NormalTok{],  }\CommentTok{\# all positive}
\NormalTok{)}

\CommentTok{\# If leakage detected → add negative point}
\NormalTok{lumen\_mask\_refined }\OperatorTok{=}\NormalTok{ model.segment(}
\NormalTok{    volume}\OperatorTok{=}\NormalTok{cta\_volume,}
\NormalTok{    points}\OperatorTok{=}\NormalTok{[ostium\_xyz, mid\_vessel\_xyz, distal\_xyz, vein\_xyz],}
\NormalTok{    labels}\OperatorTok{=}\NormalTok{[}\DecValTok{1}\NormalTok{, }\DecValTok{1}\NormalTok{, }\DecValTok{1}\NormalTok{, }\DecValTok{0}\NormalTok{],  }\CommentTok{\# last point is negative}
\NormalTok{)}
\end{Highlighting}
\end{Shaded}

\paragraph{B. Dual-Wall Nested Segmentation (Lumen →
EEM)}\label{b.-dual-wall-nested-segmentation-lumen-eem}

Critical for calculating \textbf{stenosis \%} and \textbf{plaque
burden}:

\begin{quote}
\textbf{Stenosis \% definition:} Report \textbf{diameter stenosis} = (1
− D\_min / D\_ref) × 100, consistent with CAD-RADS 2.0 grading. D\_ref
is the mean diameter of the nearest healthy proximal and distal
reference segments. Do \textbf{not} use area stenosis (more sensitive
but not the CAD-RADS standard), and document the reference segment
selection method explicitly --- QAngio CT (MEDIS) uses interpolated
reference; our pipeline must match this for Dice/correlation comparisons
against expert measurements.
\end{quote}

\begin{enumerate}
\def\labelenumi{\arabic{enumi}.}
\tightlist
\item
  \textbf{Step 1:} Segment lumen (bright contrast, easier target)
\item
  \textbf{Step 2:} Use lumen mask as \textbf{dense prompt} → expand
  outward to EEM
\item
  \textbf{Result:}
  \texttt{vessel\_wall\ =\ outer\_mask\ \&\ \textasciitilde{}lumen\_mask}
  → plaque volume
\end{enumerate}

\begin{quote}
\textbf{Implementation warning:} SAM-Med3D does \textbf{not} natively
support a \texttt{dense\_prompt} argument in its published codebase. The
pseudocode below shows the \emph{intended} design. Two implementation
paths exist: (a) custom modification of the prompt encoder to accept a
prior mask, or (b) a two-stage pipeline where the lumen mask surface is
sampled into additional positive point prompts. This is a research
contribution we must implement before the code below is functional.
\end{quote}

\begin{Shaded}
\begin{Highlighting}[]
\CommentTok{\# PSEUDOCODE — dense\_prompt is not yet implemented in SAM{-}Med3D}

\CommentTok{\# Step 1: Lumen}
\NormalTok{lumen\_mask }\OperatorTok{=}\NormalTok{ model.segment(volume, points}\OperatorTok{=}\NormalTok{[ostium, distal], labels}\OperatorTok{=}\NormalTok{[}\DecValTok{1}\NormalTok{, }\DecValTok{1}\NormalTok{])}

\CommentTok{\# Step 2: Outer wall using lumen as prior (requires custom prompt encoder)}
\NormalTok{outer\_mask }\OperatorTok{=}\NormalTok{ model.segment(}
\NormalTok{    volume,}
\NormalTok{    points}\OperatorTok{=}\NormalTok{[ostium],}
\NormalTok{    dense\_prompt}\OperatorTok{=}\NormalTok{lumen\_mask,  }\CommentTok{\# NOT natively supported — must be implemented}
\NormalTok{    labels}\OperatorTok{=}\NormalTok{[}\DecValTok{1}\NormalTok{]}
\NormalTok{)}

\CommentTok{\# Step 3: Derive vessel wall}
\NormalTok{vessel\_wall }\OperatorTok{=}\NormalTok{ outer\_mask }\OperatorTok{\&} \OperatorTok{\textasciitilde{}}\NormalTok{lumen\_mask}
\end{Highlighting}
\end{Shaded}

\paragraph{C. Post-Processing: Plaque Characterisation by HU
Thresholds}\label{c.-post-processing-plaque-characterisation-by-hu-thresholds}

\begin{Shaded}
\begin{Highlighting}[]
\KeywordTok{def}\NormalTok{ characterise\_plaque(volume, vessel\_wall\_mask):}
    \CommentTok{"""Classify plaque components by HU value within the vessel wall mask."""}
\NormalTok{    hu\_values }\OperatorTok{=}\NormalTok{ volume[vessel\_wall\_mask }\OperatorTok{\textgreater{}} \DecValTok{0}\NormalTok{]}

\NormalTok{    calcified    }\OperatorTok{=}\NormalTok{ (hu\_values }\OperatorTok{\textgreater{}} \DecValTok{130}\NormalTok{).}\BuiltInTok{sum}\NormalTok{()}
\NormalTok{    fibrous      }\OperatorTok{=}\NormalTok{ ((hu\_values }\OperatorTok{\textgreater{}=} \DecValTok{60}\NormalTok{) }\OperatorTok{\&}\NormalTok{ (hu\_values }\OperatorTok{\textless{}=} \DecValTok{130}\NormalTok{)).}\BuiltInTok{sum}\NormalTok{()}
\NormalTok{    lipid\_rich   }\OperatorTok{=}\NormalTok{ (hu\_values }\OperatorTok{\textless{}} \DecValTok{60}\NormalTok{).}\BuiltInTok{sum}\NormalTok{()}
\NormalTok{    total        }\OperatorTok{=}\NormalTok{ vessel\_wall\_mask.}\BuiltInTok{sum}\NormalTok{()}

    \ControlFlowTok{return}\NormalTok{ \{}
        \StringTok{"calcified\_pct"}\NormalTok{: calcified }\OperatorTok{/}\NormalTok{ total }\OperatorTok{*} \DecValTok{100}\NormalTok{,}
        \StringTok{"fibrous\_pct"}\NormalTok{:   fibrous }\OperatorTok{/}\NormalTok{ total }\OperatorTok{*} \DecValTok{100}\NormalTok{,}
        \StringTok{"lipid\_rich\_pct"}\NormalTok{: lipid\_rich }\OperatorTok{/}\NormalTok{ total }\OperatorTok{*} \DecValTok{100}\NormalTok{,}
        \StringTok{"high\_risk"}\NormalTok{: (lipid\_rich }\OperatorTok{/}\NormalTok{ total) }\OperatorTok{\textgreater{}} \FloatTok{0.04}\NormalTok{,  }\CommentTok{\# \textgreater{}4\% LAP = vulnerable}
\NormalTok{    \}}
\end{Highlighting}
\end{Shaded}

\begin{quote}
\textbf{HU calibration warning:} The thresholds above (\textless{} 60
LAP, 60--130 fibrous, \textgreater{} 130 calcified) are defined for
\textbf{standard 120 kVp with soft-kernel reconstruction}. DISCHARGE is
a multi-centre trial with varying tube voltages (80--140 kVp) and
reconstruction kernels (soft vs.~sharp). Sharp-kernel reconstruction
shifts apparent HU upward and increases blooming around calcification.
These thresholds require per-centre calibration against phantom
measurements or paired soft/sharp kernel scans before cross-site plaque
comparisons are valid. Record the reconstruction kernel and tube voltage
for every DISCHARGE case in the metadata.
\end{quote}

\subsubsection{DISCHARGE-Specific Processing
Considerations}\label{discharge-specific-processing-considerations}

\begin{longtable}[]{@{}
  >{\raggedright\arraybackslash}p{(\columnwidth - 2\tabcolsep) * \real{0.6522}}
  >{\raggedright\arraybackslash}p{(\columnwidth - 2\tabcolsep) * \real{0.3478}}@{}}
\toprule\noalign{}
\begin{minipage}[b]{\linewidth}\raggedright
Consideration
\end{minipage} & \begin{minipage}[b]{\linewidth}\raggedright
Detail
\end{minipage} \\
\midrule\noalign{}
\endhead
\bottomrule\noalign{}
\endlastfoot
\textbf{Scanner heterogeneity} & Multi-centre trial → varying scanner
vendors, protocols, contrast timing \\
\textbf{Reconstruction kernels} & Soft vs.~sharp kernels affect HU
accuracy for plaque \\
\textbf{Contrast timing} & Early arterial phase optimal; late phase
reduces lumen-wall contrast \\
\textbf{ECG gating} & Prospective vs.~retrospective gating affects
motion artefact severity \\
\textbf{Data format} & DICOM (clinical) → convert to NIfTI.gz for model
input \\
\textbf{Annotations} & MEDIS QAngio CT (expert contours) available for
subset → ground truth \\
\end{longtable}

\begin{center}\rule{0.5\linewidth}{0.5pt}\end{center}

\subsection{Clinical Domain B: Prostate mpMRI
Segmentation}\label{clinical-domain-b-prostate-mpmri-segmentation}

\subsubsection{Imaging Standard}\label{imaging-standard}

\textbf{Multiparametric MRI (mpMRI)} is the clinical standard for
prostate imaging, using: - \textbf{T2-weighted (T2W):} Anatomical detail
--- zonal anatomy - \textbf{Diffusion-weighted imaging (DWI) + ADC map:}
Cellularity --- lesion detection - \textbf{Dynamic contrast-enhanced
(DCE):} Vascularity --- supplementary

Reporting follows \textbf{PI-RADS v2.1} (Prostate Imaging-Reporting and
Data System).

\subsubsection{Anatomical Segmentation --- The
``Zones''}\label{anatomical-segmentation-the-zones}

The prostate is divided into distinct zones with different MRI
appearances and cancer risk:

\begin{longtable}[]{@{}
  >{\raggedright\arraybackslash}p{(\columnwidth - 8\tabcolsep) * \real{0.0952}}
  >{\raggedright\arraybackslash}p{(\columnwidth - 8\tabcolsep) * \real{0.2063}}
  >{\raggedright\arraybackslash}p{(\columnwidth - 8\tabcolsep) * \real{0.2063}}
  >{\raggedright\arraybackslash}p{(\columnwidth - 8\tabcolsep) * \real{0.2540}}
  >{\raggedright\arraybackslash}p{(\columnwidth - 8\tabcolsep) * \real{0.2381}}@{}}
\toprule\noalign{}
\begin{minipage}[b]{\linewidth}\raggedright
Zone
\end{minipage} & \begin{minipage}[b]{\linewidth}\raggedright
Abbreviation
\end{minipage} & \begin{minipage}[b]{\linewidth}\raggedright
Cancer Risk
\end{minipage} & \begin{minipage}[b]{\linewidth}\raggedright
T2W Appearance
\end{minipage} & \begin{minipage}[b]{\linewidth}\raggedright
Clinical Role
\end{minipage} \\
\midrule\noalign{}
\endhead
\bottomrule\noalign{}
\endlastfoot
\textbf{Peripheral Zone} & PZ & 70--75 \% of cancers & Bright (high
signal) & Primary cancer surveillance region \\
\textbf{Transition Zone} & TZ & 20--25 \% of cancers & Heterogeneous
(BPH nodules) & BPH assessment, cancer in older men \\
\textbf{Central Zone} & CZ & \textless{} 5 \% of cancers & Low signal
(dense stroma) & Indistinguishable from TZ on MRI --- grouped with TZ
per PI-RADS v2.1 \\
\textbf{Anterior Fibromuscular Stroma} & AFMS & Non-glandular & Very low
signal & Can be invaded by anterior tumours \\
\end{longtable}

\subsubsection{Pathology Segmentation --- The
``Lesions''}\label{pathology-segmentation-the-lesions}

When segmenting pathology, the target is \textbf{clinically significant
prostate cancer (csPCa)}:

\begin{longtable}[]{@{}
  >{\raggedright\arraybackslash}p{(\columnwidth - 4\tabcolsep) * \real{0.2708}}
  >{\raggedright\arraybackslash}p{(\columnwidth - 4\tabcolsep) * \real{0.2708}}
  >{\raggedright\arraybackslash}p{(\columnwidth - 4\tabcolsep) * \real{0.4583}}@{}}
\toprule\noalign{}
\begin{minipage}[b]{\linewidth}\raggedright
Lesion Type
\end{minipage} & \begin{minipage}[b]{\linewidth}\raggedright
Description
\end{minipage} & \begin{minipage}[b]{\linewidth}\raggedright
Clinical Significance
\end{minipage} \\
\midrule\noalign{}
\endhead
\bottomrule\noalign{}
\endlastfoot
\textbf{Index Lesion} & Largest / most aggressive tumour & Primary
target for biopsy and treatment \\
\textbf{Satellite Lesions} & Secondary foci (prostate cancer is often
multifocal) & May affect treatment strategy \\
\textbf{Extracapsular Extension (ECE)} & Tumour breaches the prostatic
capsule & Staging: T3a --- affects surgical planning \\
\textbf{Seminal Vesicle Invasion (SVI)} & Tumour extends into seminal
vesicles & Staging: T3b --- impacts prognosis \\
\end{longtable}

\subsubsection{Organs at Risk (OARs) for
Radiotherapy}\label{organs-at-risk-oars-for-radiotherapy}

\begin{longtable}[]{@{}
  >{\raggedright\arraybackslash}p{(\columnwidth - 4\tabcolsep) * \real{0.2973}}
  >{\raggedright\arraybackslash}p{(\columnwidth - 4\tabcolsep) * \real{0.3514}}
  >{\raggedright\arraybackslash}p{(\columnwidth - 4\tabcolsep) * \real{0.3514}}@{}}
\toprule\noalign{}
\begin{minipage}[b]{\linewidth}\raggedright
Structure
\end{minipage} & \begin{minipage}[b]{\linewidth}\raggedright
Abbreviation
\end{minipage} & \begin{minipage}[b]{\linewidth}\raggedright
Why Segment?
\end{minipage} \\
\midrule\noalign{}
\endhead
\bottomrule\noalign{}
\endlastfoot
\textbf{Neurovascular Bundles} & NVB & Nerve-sparing surgery ---
preserve potency \\
\textbf{Rectal Wall} & Rectum\_Wall & Monitor tumour--rectum distance \\
\textbf{Bladder Neck} & Bladder\_Neck & Preserve urinary continence \\
\end{longtable}

\subsubsection{Segmentation Class Labels for the AI
Model}\label{segmentation-class-labels-for-the-ai-model}

\begin{longtable}[]{@{}
  >{\raggedright\arraybackslash}p{(\columnwidth - 6\tabcolsep) * \real{0.2600}}
  >{\raggedright\arraybackslash}p{(\columnwidth - 6\tabcolsep) * \real{0.2400}}
  >{\raggedright\arraybackslash}p{(\columnwidth - 6\tabcolsep) * \real{0.2000}}
  >{\raggedright\arraybackslash}p{(\columnwidth - 6\tabcolsep) * \real{0.3000}}@{}}
\toprule\noalign{}
\begin{minipage}[b]{\linewidth}\raggedright
Segment Name
\end{minipage} & \begin{minipage}[b]{\linewidth}\raggedright
Class Label
\end{minipage} & \begin{minipage}[b]{\linewidth}\raggedright
Modality
\end{minipage} & \begin{minipage}[b]{\linewidth}\raggedright
Clinical Goal
\end{minipage} \\
\midrule\noalign{}
\endhead
\bottomrule\noalign{}
\endlastfoot
Whole Gland & \texttt{Prostate\_WG} & T2W & Volume / PSA density
calculation \\
Peripheral Zone & \texttt{PZ} & T2W & Cancer surveillance background \\
Transition Zone & \texttt{TZ} & T2W & BPH assessment background \\
Suspicious Lesion &
\texttt{Lesion\_PIRADS\_\{3\textbackslash{}\textbar{}4\textbackslash{}\textbar{}5\}}
& DWI/ADC & Targeted biopsy (MR-US fusion) \\
Seminal Vesicles & \texttt{SV} & T2W & Local staging (T3b) \\
Neurovascular Bundles & \texttt{NVB} & T2W & Nerve-sparing planning \\
Rectal Wall & \texttt{OAR\_Rectum} & T2W & Radiotherapy constraints \\
\end{longtable}

\subsubsection{SAM-Med3D-turbo Prompting Strategy for Prostate
mpMRI}\label{sam-med3d-turbo-prompting-strategy-for-prostate-mpmri}

\textbf{Context-dependent prompting:} The same model must switch
behaviour based on \emph{where} the user clicks and \emph{which
sequence} is active.

\paragraph{Scenario 1: Anatomical Zone Segmentation
(T2W)}\label{scenario-1-anatomical-zone-segmentation-t2w}

\begin{verbatim}
1. User clicks bright outer rim on T2W axial
   → Model returns PZ mask
   → Expected Dice: ~0.90 (large, high-contrast structure)

2. User clicks central heterogeneous region on T2W
   → Model returns TZ mask

3. If mask leaks into rectum → place negative point on rectum
\end{verbatim}

\paragraph{Scenario 2: Lesion Segmentation
(DWI/ADC)}\label{scenario-2-lesion-segmentation-dwiadc}

\begin{quote}
\textbf{Implementation warning:} Step 2 below references using a prior
PZ mask as a dense prompt. This is the same unimplemented
\texttt{dense\_prompt} issue described in §3.4.2 --- the same two
implementation paths apply: (a) custom prompt encoder modification, or
(b) sampling the PZ mask boundary into additional point prompts.
\end{quote}

\begin{verbatim}
1. User clicks hypointense spot on ADC map
   → Model returns lesion mask (PI-RADS ≥4 region)

2. [PLANNED] Use prior PZ mask as dense prompt context
   → Constrains lesion to within prostate boundary
   → Requires dense_prompt implementation (see §3.4.2)

3. Measure lesion volume → maps to PI-RADS size criterion
\end{verbatim}

\begin{quote}
\textbf{Critical note:} SAM-Med3D was pre-trained on \textbf{MR data}
(SA-Med3D-140K includes MR modalities). However, the dataset card does
not specify which MR sequences or whether prostate mpMRI is represented.
Zero-shot performance on prostate zones may be moderate; fine-tuning on
Charité prostate data will likely be necessary. Whole-gland segmentation
(a large, well-defined structure) should work well zero-shot; zonal
segmentation (PZ vs.~TZ) is harder due to subtle signal differences.
\end{quote}

\subsubsection{Prostate vs.~Coronary: Difficulty
Comparison}\label{prostate-vs.-coronary-difficulty-comparison}

\begin{longtable}[]{@{}
  >{\raggedright\arraybackslash}p{(\columnwidth - 4\tabcolsep) * \real{0.2857}}
  >{\raggedright\arraybackslash}p{(\columnwidth - 4\tabcolsep) * \real{0.3571}}
  >{\raggedright\arraybackslash}p{(\columnwidth - 4\tabcolsep) * \real{0.3571}}@{}}
\toprule\noalign{}
\begin{minipage}[b]{\linewidth}\raggedright
Factor
\end{minipage} & \begin{minipage}[b]{\linewidth}\raggedright
Prostate
\end{minipage} & \begin{minipage}[b]{\linewidth}\raggedright
Coronary
\end{minipage} \\
\midrule\noalign{}
\endhead
\bottomrule\noalign{}
\endlastfoot
Structure size & 20--80 mL (large) & 1.5--4 mm diameter (tiny) \\
Motion & None (static pelvis) & Cardiac motion (60--80 bpm) \\
Contrast & Good (gland vs.~fat) & Variable (plaque vs.~lumen) \\
Modality & MRI (multi-sequence) & CT (single phase) \\
Topology & Compact, roughly ellipsoidal & Thin, tortuous, branching
tubes \\
\textbf{Expected Dice} & \textbf{\textgreater{} 0.90 (whole gland)} &
\textbf{0.75--0.85 (lumen)} \\
\end{longtable}

\begin{center}\rule{0.5\linewidth}{0.5pt}\end{center}

\subsection{Technical Challenges \& Model-Aware
Solutions}\label{technical-challenges-model-aware-solutions}

\subsubsection{Challenge 1: Coronary Artery Motion
Artefacts}\label{challenge-1-coronary-artery-motion-artefacts}

\textbf{Problem:} Residual cardiac motion blurs vessel edges despite ECG
gating. RCA most affected.

\textbf{Solution:}

\begin{Shaded}
\begin{Highlighting}[]
\KeywordTok{def}\NormalTok{ motion\_robust\_segmentation(volume\_4d, heart\_rate, prompt\_points, prompt\_labels):}
    \ControlFlowTok{if}\NormalTok{ heart\_rate }\OperatorTok{\textgreater{}} \DecValTok{70}\NormalTok{:}
        \CommentTok{\# Multi{-}phase reconstruction}
\NormalTok{        phases }\OperatorTok{=}\NormalTok{ extract\_cardiac\_phases(volume\_4d, num\_phases}\OperatorTok{=}\DecValTok{10}\NormalTok{)}
\NormalTok{        masks }\OperatorTok{=}\NormalTok{ [model.segment(phase, points}\OperatorTok{=}\NormalTok{prompt\_points, labels}\OperatorTok{=}\NormalTok{prompt\_labels) }\ControlFlowTok{for}\NormalTok{ phase }\KeywordTok{in}\NormalTok{ phases]}
        \CommentTok{\# Temporal median filter removes motion ghosts}
        \ControlFlowTok{return}\NormalTok{ np.median(masks, axis}\OperatorTok{=}\DecValTok{0}\NormalTok{) }\OperatorTok{\textgreater{}} \FloatTok{0.5}
    \ControlFlowTok{else}\NormalTok{:}
        \ControlFlowTok{return}\NormalTok{ model.segment(volume\_4d[:,:,:,}\DecValTok{0}\NormalTok{], points}\OperatorTok{=}\NormalTok{prompt\_points, labels}\OperatorTok{=}\NormalTok{prompt\_labels)}
\end{Highlighting}
\end{Shaded}

\begin{quote}
\textbf{Topology warning:} Pixel-wise median of binary masks across 10
phases is \textbf{not} topologically safe for thin structures. A distal
coronary voxel present in only 4/10 phases will be excluded by the
median, potentially severing the centreline. Preferred alternative:
select the optimal cardiac phase (75 \% R-R for RCA, 65 \% R-R for
LAD/LCx at HR \textless{} 70) rather than fusing all phases. If
multi-phase fusion is required, apply connected-component analysis
post-median and re-link broken segments by dilation along the
centreline.
\end{quote}

\textbf{Additional strategies:} - Edge-preserving denoising (bilateral
filter) pre-processing - Multi-point prompts every 5--10 mm along vessel
to guide through corrupted regions - Auto-flag cases with edge sharpness
\textless{} threshold for expert review

\subsubsection{Challenge 2: Low-Attenuation Plaque
Detection}\label{challenge-2-low-attenuation-plaque-detection}

\textbf{Problem:} Lipid-rich plaque (20--50 HU) nearly invisible against
myocardium (50--70 HU).

\textbf{Solution:} Multi-stage approach: 1. Segment lumen and outer wall
first (high-contrast boundaries) 2. Apply HU thresholding \emph{within}
the vessel wall mask 3. Use SAM-Med3D's ViT features for texture-based
refinement (fine-tuning required)

\subsubsection{Challenge 3: Dual-Wall
Segmentation}\label{challenge-3-dual-wall-segmentation}

\textbf{Problem:} Must segment both lumen and outer wall; wall thickness
only 0.5--3 mm (1--5 voxels).

\textbf{Solution:} Sequential prompting (see §3.4.2 above).

\subsubsection{Challenge 4: Prostate Zone
Boundaries}\label{challenge-4-prostate-zone-boundaries}

\textbf{Problem:} PZ-TZ boundary is a gradual signal transition, not a
sharp edge.

\textbf{Solution:} - Train multi-class model on annotated Charité
prostate data - Use morphological priors (PZ wraps around TZ
inferolaterally) - Negative prompts at zone transitions to sharpen
boundaries

\subsubsection{Challenge 5: Model Limitations --- Honest
Assessment}\label{challenge-5-model-limitations-honest-assessment}

\begin{longtable}[]{@{}
  >{\raggedright\arraybackslash}p{(\columnwidth - 4\tabcolsep) * \real{0.3548}}
  >{\raggedright\arraybackslash}p{(\columnwidth - 4\tabcolsep) * \real{0.2581}}
  >{\raggedright\arraybackslash}p{(\columnwidth - 4\tabcolsep) * \real{0.3871}}@{}}
\toprule\noalign{}
\begin{minipage}[b]{\linewidth}\raggedright
Limitation
\end{minipage} & \begin{minipage}[b]{\linewidth}\raggedright
Impact
\end{minipage} & \begin{minipage}[b]{\linewidth}\raggedright
Mitigation
\end{minipage} \\
\midrule\noalign{}
\endhead
\bottomrule\noalign{}
\endlastfoot
No coronary CTA in pre-training & Zero-shot may underperform & Fine-tune
on DISCHARGE annotations \\
Binary mask output only & Cannot directly predict PI-RADS score &
Post-processing pipeline (volume, ADC stats) \\
No \texttt{dense\_prompt} in codebase & Lumen-as-prior strategy needs
engineering & Custom prompt encoder modification \\
128³ patch size constraint & Coronary arteries span \textgreater{} 128
voxels (300--400 voxels at 0.5 mm) & Sliding window with 50 \% overlap;
Gaussian-weighted blending at boundaries to prevent double-thick seams;
connected-component check post-stitch \\
No multi-class output & One mask per inference call & Multiple
sequential inferences per case \\
\end{longtable}

\begin{center}\rule{0.5\linewidth}{0.5pt}\end{center}

\subsection{Clinical Workflows}\label{clinical-workflows}

\subsubsection{Workflow 1: Interactive Coronary Segmentation
(Radiologist)}\label{workflow-1-interactive-coronary-segmentation-radiologist}

\begin{enumerate}
\def\labelenumi{\arabic{enumi}.}
\tightlist
\item
  Radiologist opens browser → logs in (Charité SSO)
\item
  Loads DISCHARGE CCTA scan from PACS
\item
  Clicks proximal LAD + distal tip → AI segments entire lumen in
  \textless{} 2 s
\item
  Optionally: places negative point to prune vein leakage
\item
  Clicks ``Expand to Wall'' → second inference → outer wall mask
  \textbf{{[}PLANNED --- requires \texttt{dense\_prompt} implementation;
  see §3.4.2{]}}
\item
  Right panel shows: stenosis \% (diameter stenosis, CAD-RADS 2.0),
  plaque composition, volume
\item
  Exports segmentation (NIfTI / DICOM-SEG / CSV report)
\end{enumerate}

\subsubsection{Workflow 2: Batch Processing for DISCHARGE
Research}\label{workflow-2-batch-processing-for-discharge-research}

\begin{enumerate}
\def\labelenumi{\arabic{enumi}.}
\tightlist
\item
  Research coordinator uploads 100 DISCHARGE cases
\item
  \textbf{Automated ostium seeding:} template-based registration to a
  reference coronary atlas provides initial prompt coordinates; cases
  where registration confidence is low are queued for manual prompt
  review before inference
\item
  AI processes overnight (Celery + multi-GPU batch mode)
\item
  Quality control: auto-flag cases with \textbf{disconnected mask
  components, mask volume outside 3σ of cohort distribution, or model
  confidence score below threshold} (ground truth is not available in
  batch mode --- Dice cannot be computed)
\item
  Expert reviews flagged cases → corrections exported as corrected
  per-structure binary NIfTI masks → feed active-learning loop
\item
  Export refined segmentations for MACE prediction analysis
\end{enumerate}

\subsubsection{Workflow 3: MEDIS TXT + Mesh + Straightened MPR
(Reference
Contours)}\label{workflow-3-medis-txt-mesh-straightened-mpr-reference-contours}

\begin{enumerate}
\def\labelenumi{\arabic{enumi}.}
\tightlist
\item
  Load CCTA volume (NIfTI.gz) in browser
\item
  Load MEDIS TXT file (expert contour rings)
\item
  Client-side mesh generation: parse TXT → NVMesh (50--100 ms)
\item
  Overlay lumen + vessel wall meshes on CCTA
\item
  Generate straightened MPR for longitudinal plaque assessment
\end{enumerate}

\begin{center}\rule{0.5\linewidth}{0.5pt}\end{center}

\subsection{System Architecture}\label{system-architecture}

\subsubsection{High-Level Overview}\label{high-level-overview}

\begin{verbatim}
┌─────────────────────────────────────────────────────────────┐
│                    Charité Browser (HTTPS)                   │
│  ┌───────────────────────────────────────────────────────┐  │
│  │  Frontend: TypeScript + Vite + Niivue 0.66            │  │
│  │  - Volume rendering (WebGL2, 60 FPS)                  │  │
│  │  - Interactive prompting (click → 3D coordinates)     │  │
│  │  - Mesh overlay, quad-view MPR                        │  │
│  └───────────────┬───────────────────────────────────────┘  │
└──────────────────┼──────────────────────────────────────────┘
                   │ REST API (JSON + NIfTI blobs)
┌──────────────────┼──────────────────────────────────────────┐
│  Backend: FastAPI + PyTorch (on-premise GPU cluster)        │
│  ┌───────────────┼───────────────────────────────────────┐  │
│  │  API Gateway (auth, rate-limit, CORS)                 │  │
│  ├───────────────┼───────────────────────────────────────┤  │
│  │  SAM-Med3D-turbo  │  nnU-Net (prior)  │  HU Pipeline │  │
│  ├───────────────┼───────────────────────────────────────┤  │
│  │  Redis (embedding cache) │ Celery (batch queue)       │  │
│  ├───────────────┼───────────────────────────────────────┤  │
│  │  NVIDIA A100 GPUs (4×)  │  DICOM/NIfTI storage       │  │
│  └───────────────────────────────────────────────────────┘  │
└─────────────────────────────────────────────────────────────┘
\end{verbatim}

\subsubsection{Frontend Directory
Structure}\label{frontend-directory-structure}

\begin{verbatim}
flow-segment-frontend/
├── src/
│   ├── main.ts                    # Entry point
│   ├── App.ts                     # Main app component
│   ├── components/
│   │   ├── NiivueViewer.ts        # Niivue canvas wrapper (single/quad)
│   │   ├── Toolbar.ts             # Top toolbar
│   │   ├── SegmentPanel.ts        # AI segmentation controls
│   │   ├── ResultsPanel.ts        # Plaque analysis / zone selector
│   │   ├── PromptHistory.ts       # Point prompt log + undo
│   │   └── MPRView.ts             # Multi-planar reconstruction
│   ├── services/
│   │   ├── api.ts                 # Backend API client
│   │   ├── niivue.ts              # Niivue init & config
│   │   ├── loader.ts              # NIfTI.gz + DICOM loading
│   │   ├── medisParser.ts         # MEDIS TXT parsing
│   │   ├── medisMeshDirect.ts     # Client-side contour→mesh
│   │   ├── straightenedMPR.ts     # CPR (see §12)
│   │   └── auth.ts                # LDAP/SSO
│   ├── types/
│   │   ├── volume.ts
│   │   ├── segmentation.ts
│   │   ├── mesh.ts
│   │   └── api.ts
│   └── utils/
│       ├── coordinates.ts         # 3D coordinate transforms
│       ├── meshGenerator.ts       # Marching cubes (vtk.js)
│       └── export.ts              # NIfTI / DICOM-SEG / CSV export
├── package.json
├── tsconfig.json
├── vite.config.ts
└── index.html
\end{verbatim}

\subsubsection{Backend Directory
Structure}\label{backend-directory-structure}

\begin{verbatim}
flow-segment-backend/
├── app/
│   ├── main.py                    # FastAPI application
│   ├── config.py                  # Environment config
│   ├── models/
│   │   ├── sam_med3d.py           # SAM-Med3D-turbo wrapper
│   │   ├── nnu_net.py             # nnU-Net (anatomical prior)
│   │   └── plaque_analyser.py     # HU-based plaque classification
│   ├── api/
│   │   ├── segment.py             # Segmentation endpoints
│   │   ├── volumes.py             # Volume management
│   │   ├── mesh.py                # Mesh generation (nii2mesh)
│   │   └── auth.py                # LDAP/SSO
│   └── services/
│       ├── cache.py               # Redis embedding cache
│       ├── dicom_processor.py     # DICOM → NIfTI (SimpleITK)
│       ├── mesh_generator.py      # nii2mesh wrapper
│       └── registration.py        # Elastix (longitudinal)
├── Dockerfile
├── docker-compose.yml
├── requirements.txt
└── README.md
\end{verbatim}

\begin{center}\rule{0.5\linewidth}{0.5pt}\end{center}

\subsection{\texorpdfstring{Frontend: Niivue Viewer \&
UI/UX\textsuperscript{5}}{Frontend: Niivue Viewer \& UI/UX5}}\label{frontend-niivue-viewer-uiux-niivue2024}

\subsubsection{Key Niivue v0.66.0
Capabilities}\label{key-niivue-v0.66.0-capabilities}

\begin{itemize}
\tightlist
\item
  WebGL2 rendering: 60 FPS for 512³ volumes
\item
  \texttt{onLocationChange} / \texttt{onMouseUp} → capture 3D mm
  coordinates for prompts
\item
  Multi-planar reconstruction (axial / coronal / sagittal sync)
\item
  \texttt{nv.addMesh()} for 3D surface overlay with adjustable opacity
\item
  Full TypeScript definitions
\item
  In-browser DICOM via dcm2niix-wasm
\end{itemize}

\subsubsection{UI Layout: Single View
(Default)}\label{ui-layout-single-view-default}

\begin{verbatim}
┌─────────────────────────────────────────────────────────────┐
│  [Logo] Segment    [Load NII/DCM] [Save] [Settings] [User] [◧ Quad]│
├─────────────────────────────────────────────────────────────┤
│  ┌─────────────────────────────────┐  ┌───────────────────┐ │
│  │                                 │  │  Segmentation     │ │
│  │   Niivue 3D Canvas             │  │  ┌──────────────┐  │ │
│  │   [Click to place prompt]      │  │  │ ● Point      │  │ │
│  │                                 │  │  │ □ Box        │  │ │
│  │                                 │  │  │ ▶ Segment    │  │ │
│  │                                 │  │  └──────────────┘  │ │
│  │                                 │  │                    │ │
│  │                                 │  │  Prompt History    │ │
│  │                                 │  │  + (120.5, 85, 42) │ │
│  │                                 │  │  - (118.2, 90, 42) │ │
│  │                                 │  │                    │ │
│  │                                 │  │  Results           │ │
│  │                                 │  │  Stenosis: 62 %    │ │
│  │                                 │  │  Calc: 45 %        │ │
│  │                                 │  │  LAP: 18 %         │ │
│  └─────────────────────────────────┘  └───────────────────┘ │
└─────────────────────────────────────────────────────────────┘
\end{verbatim}

\subsubsection{UI Layout: Quad View (MPR +
3D)}\label{ui-layout-quad-view-mpr-3d}

\begin{verbatim}
┌─────────────────────────────────────────────────────────────┐
│  [Logo] Segment    [Load NII/DCM] [Save] [Settings] [User] [◫ 1x1]│
├─────────────────────────────────────────────────────────────┤
│  ┌──────────────────┐ ┌──────────────────┐  ┌────────────┐ │
│  │  Axial           │ │  Sagittal        │  │ Controls   │ │
│  │  + crosshair     │ │  + crosshair     │  │            │ │
│  ├──────────────────┤ ├──────────────────┤  │ Zone/Vessel│ │
│  │  Coronal         │ │  3D Render       │  │ Selector   │ │
│  │  + crosshair     │ │  + mesh overlay  │  │            │ │
│  └──────────────────┘ └──────────────────┘  └────────────┘ │
└─────────────────────────────────────────────────────────────┘
\end{verbatim}

\subsubsection{Design Principles}\label{design-principles}

\begin{itemize}
\tightlist
\item
  \textbf{Dark theme} (\#1a1a1a) --- reduces eye strain for long
  sessions
\item
  \textbf{Minimal chrome} --- maximise canvas area
\item
  \textbf{Radiologist-first} --- optimised for clinical workflow
\item
  Primary: Blue (\#3b82f6), Accent: Green (\#10b981), Warning: Orange
  (\#f59e0b), Error: Red (\#ef4444)
\end{itemize}

\subsubsection{Keyboard Shortcuts}\label{keyboard-shortcuts}

\begin{longtable}[]{@{}ll@{}}
\toprule\noalign{}
Key & Action \\
\midrule\noalign{}
\endhead
\bottomrule\noalign{}
\endlastfoot
\texttt{P} & Place positive prompt point \\
\texttt{N} & Place negative prompt point \\
\texttt{Enter} & Run segmentation with current prompts \\
\texttt{Z} & Undo last prompt point \\
\texttt{Escape} & Cancel current prompt session \\
\texttt{Q} & Toggle single / quad view \\
\texttt{W} / \texttt{L} & Adjust window / level (hold + drag) \\
\texttt{Space} & Toggle mesh overlay visibility \\
\end{longtable}

\subsubsection{General-Purpose Rotatable Volume
Viewer}\label{general-purpose-rotatable-volume-viewer}

Port of legacy \texttt{viewer.py} (PyQtGraph) to WebGL:

\begin{itemize}
\tightlist
\item
  Free rotation via mouse drag (Rodrigues' rotation formula)
\item
  Three-panel orthogonal views (YZ, XZ, XY) --- all rotate together
\item
  Drag modes: Rotation (0), Paint/Label (1), Window/Level (2)
\item
  Real-time volume resampling (SimpleITK → WebGL equivalent)
\item
  Target: 30--60 FPS during rotation
\end{itemize}

\begin{center}\rule{0.5\linewidth}{0.5pt}\end{center}

\subsection{Backend: Inference Pipeline \&
API}\label{backend-inference-pipeline-api}

\subsubsection{API Endpoints}\label{api-endpoints}

\begin{verbatim}
# Segmentation (AI)
POST /api/segment/point        # Point-based prompting (SAM-Med3D)
POST /api/segment/box          # Bounding box prompting
POST /api/segment/refine       # Refine with negative points

# Volume Management
POST /api/volumes/upload       # Upload DICOM / NIfTI
GET  /api/volumes/{id}         # Retrieve processed volume

# Mesh Generation
POST /api/mesh/from-mask       # Mask → MZ3 mesh (nii2mesh)
POST /api/mesh/quick           # Fast preview (marching cubes)

# MEDIS TXT Processing
POST /api/medis/upload         # Upload MEDIS TXT
POST /api/medis/to-mesh        # Convert contours to mesh

# Straightened MPR
POST /api/straighten/create    # Centerline → straightened volume

# Batch Processing
POST /api/batch/process        # Queue batch of cases
GET  /api/batch/{job_id}       # Job status

# Authentication
POST /api/auth/login           # LDAP / Charité SSO
GET  /api/health               # Health check + GPU status
\end{verbatim}

\subsubsection{Segmentation Endpoint
(Core)}\label{segmentation-endpoint-core}

\begin{Shaded}
\begin{Highlighting}[]
\ImportTok{from}\NormalTok{ typing }\ImportTok{import}\NormalTok{ Optional}
\ImportTok{from}\NormalTok{ fastapi }\ImportTok{import}\NormalTok{ APIRouter}
\ImportTok{from}\NormalTok{ pydantic }\ImportTok{import}\NormalTok{ BaseModel}
\ImportTok{import}\NormalTok{ nibabel }\ImportTok{as}\NormalTok{ nib}
\ImportTok{import}\NormalTok{ numpy }\ImportTok{as}\NormalTok{ np}
\ImportTok{import}\NormalTok{ medim}

\NormalTok{router }\OperatorTok{=}\NormalTok{ APIRouter()}

\KeywordTok{class}\NormalTok{ SegmentRequest(BaseModel):}
\NormalTok{    volume\_id: }\BuiltInTok{str}
\NormalTok{    coordinates: }\BuiltInTok{list}\NormalTok{[}\BuiltInTok{list}\NormalTok{[}\BuiltInTok{float}\NormalTok{]]   }\CommentTok{\# [[x,y,z], ...] in mm — one or more points}
\NormalTok{    labels: Optional[}\BuiltInTok{list}\NormalTok{[}\BuiltInTok{int}\NormalTok{]] }\OperatorTok{=} \VariableTok{None}  \CommentTok{\# 1=positive, 0=negative per point; defaults to all{-}1}

\CommentTok{\# Model loaded once at startup}
\NormalTok{model }\OperatorTok{=}\NormalTok{ medim.create\_model(}
    \StringTok{"SAM{-}Med3D"}\NormalTok{, pretrained}\OperatorTok{=}\VariableTok{True}\NormalTok{,}
\NormalTok{    checkpoint\_path}\OperatorTok{=}\StringTok{"app/models/checkpoints/sam\_med3d\_turbo.pth"}
\NormalTok{)}
\NormalTok{model }\OperatorTok{=}\NormalTok{ model.cuda().half().}\BuiltInTok{eval}\NormalTok{()}

\AttributeTok{@router.post}\NormalTok{(}\StringTok{"/api/segment/point"}\NormalTok{)}
\ControlFlowTok{async} \KeywordTok{def}\NormalTok{ segment\_point(req: SegmentRequest):}
    \CommentTok{\# Load volume from cache/disk}
\NormalTok{    vol\_nii }\OperatorTok{=}\NormalTok{ nib.load(}\SpecialStringTok{f"/data/volumes/}\SpecialCharTok{\{}\NormalTok{req}\SpecialCharTok{.}\NormalTok{volume\_id}\SpecialCharTok{\}}\SpecialStringTok{.nii.gz"}\NormalTok{)}
\NormalTok{    volume }\OperatorTok{=}\NormalTok{ vol\_nii.get\_fdata()}

    \CommentTok{\# Convert mm → voxel coordinates for each point}
\NormalTok{    inv\_affine }\OperatorTok{=}\NormalTok{ np.linalg.inv(vol\_nii.affine)}
\NormalTok{    voxel\_points }\OperatorTok{=}\NormalTok{ [(inv\_affine }\OperatorTok{@}\NormalTok{ [}\OperatorTok{*}\NormalTok{pt, }\DecValTok{1}\NormalTok{])[:}\DecValTok{3}\NormalTok{] }\ControlFlowTok{for}\NormalTok{ pt }\KeywordTok{in}\NormalTok{ req.coordinates]}

\NormalTok{    effective\_labels }\OperatorTok{=}\NormalTok{ req.labels }\KeywordTok{or}\NormalTok{ [}\DecValTok{1}\NormalTok{] }\OperatorTok{*} \BuiltInTok{len}\NormalTok{(voxel\_points)}

    \CommentTok{\# Run SAM{-}Med3D inference}
\NormalTok{    mask }\OperatorTok{=}\NormalTok{ model.segment(volume, points}\OperatorTok{=}\NormalTok{voxel\_points, labels}\OperatorTok{=}\NormalTok{effective\_labels)}

    \CommentTok{\# Return mask as NIfTI}
\NormalTok{    mask\_nii }\OperatorTok{=}\NormalTok{ nib.Nifti1Image(mask.astype(np.uint8), vol\_nii.affine)}
    \CommentTok{\# ... serialize and return}
\end{Highlighting}
\end{Shaded}

\begin{center}\rule{0.5\linewidth}{0.5pt}\end{center}

\subsection{MEDIS TXT Reference
Pipeline}\label{medis-txt-reference-pipeline}

\subsubsection{Format Description}\label{format-description}

MEDIS TXT contains vessel wall contours from MEDIS QAngio CT: -
\textbf{Lumen} contours: define inner wall (blood pool boundary) -
\textbf{VesselWall} contours: define outer wall (including plaque) -
Each contour: group label, slice distance, N points in 3D mm coordinates

\subsubsection{Parser (TypeScript)}\label{parser-typescript}

\begin{Shaded}
\begin{Highlighting}[]
\ImportTok{export} \KeywordTok{interface}\NormalTok{ MedisContour \{}
\NormalTok{  group}\OperatorTok{:} \StringTok{"Lumen"} \OperatorTok{|} \StringTok{"VesselWall"}\OperatorTok{;}
\NormalTok{  sliceDistance}\OperatorTok{:} \DataTypeTok{number}\OperatorTok{;}           \CommentTok{// mm along vessel}
\NormalTok{  points}\OperatorTok{:}\NormalTok{ \{ x}\OperatorTok{:} \DataTypeTok{number}\OperatorTok{;}\NormalTok{ y}\OperatorTok{:} \DataTypeTok{number}\OperatorTok{;}\NormalTok{ z}\OperatorTok{:} \DataTypeTok{number}\NormalTok{ \}[]}\OperatorTok{;}
\NormalTok{\}}

\ImportTok{export} \KeywordTok{function} \FunctionTok{parseMedisTxt}\NormalTok{(content}\OperatorTok{:} \DataTypeTok{string}\NormalTok{)}\OperatorTok{:}\NormalTok{ MedisContour[] \{}
  \KeywordTok{const}\NormalTok{ lines }\OperatorTok{=}\NormalTok{ content}\OperatorTok{.}\FunctionTok{split}\NormalTok{(}\StringTok{"}\SpecialCharTok{\textbackslash{}n}\StringTok{"}\NormalTok{)}\OperatorTok{;}
  \KeywordTok{const}\NormalTok{ contours}\OperatorTok{:}\NormalTok{ MedisContour[] }\OperatorTok{=}\NormalTok{ []}\OperatorTok{;}
  \KeywordTok{let}\NormalTok{ current}\OperatorTok{:} \BuiltInTok{Partial}\OperatorTok{\textless{}}\NormalTok{MedisContour}\OperatorTok{\textgreater{}} \OperatorTok{=}\NormalTok{ \{\}}\OperatorTok{;}
  \KeywordTok{let}\NormalTok{ points}\OperatorTok{:}\NormalTok{ \{ x}\OperatorTok{:} \DataTypeTok{number}\OperatorTok{;}\NormalTok{ y}\OperatorTok{:} \DataTypeTok{number}\OperatorTok{;}\NormalTok{ z}\OperatorTok{:} \DataTypeTok{number}\NormalTok{ \}[] }\OperatorTok{=}\NormalTok{ []}\OperatorTok{;}

  \ControlFlowTok{for}\NormalTok{ (}\KeywordTok{const}\NormalTok{ line }\KeywordTok{of}\NormalTok{ lines) \{}
    \ControlFlowTok{if}\NormalTok{ (line}\OperatorTok{.}\FunctionTok{startsWith}\NormalTok{(}\StringTok{"\# group:"}\NormalTok{)) \{}
\NormalTok{      current}\OperatorTok{.}\AttributeTok{group} \OperatorTok{=}\NormalTok{ line}\OperatorTok{.}\FunctionTok{split}\NormalTok{(}\StringTok{":"}\NormalTok{)[}\DecValTok{1}\NormalTok{]}\OperatorTok{.}\FunctionTok{trim}\NormalTok{() }\ImportTok{as} \StringTok{"Lumen"} \OperatorTok{|} \StringTok{"VesselWall"}\OperatorTok{;}
\NormalTok{    \} }\ControlFlowTok{else} \ControlFlowTok{if}\NormalTok{ (line}\OperatorTok{.}\FunctionTok{startsWith}\NormalTok{(}\StringTok{"\# SliceDistance:"}\NormalTok{)) \{}
\NormalTok{      current}\OperatorTok{.}\AttributeTok{sliceDistance} \OperatorTok{=} \PreprocessorTok{parseFloat}\NormalTok{(line}\OperatorTok{.}\FunctionTok{split}\NormalTok{(}\StringTok{":"}\NormalTok{)[}\DecValTok{1}\NormalTok{])}\OperatorTok{;}
\NormalTok{    \} }\ControlFlowTok{else} \ControlFlowTok{if}\NormalTok{ (line}\OperatorTok{.}\FunctionTok{startsWith}\NormalTok{(}\StringTok{"\# Contour index:"}\NormalTok{)) \{}
      \ControlFlowTok{if}\NormalTok{ (current}\OperatorTok{.}\AttributeTok{group} \OperatorTok{\&\&}\NormalTok{ points}\OperatorTok{.}\AttributeTok{length} \OperatorTok{\textgreater{}} \DecValTok{0}\NormalTok{) \{}
\NormalTok{        contours}\OperatorTok{.}\FunctionTok{push}\NormalTok{(\{ }\OperatorTok{...}\NormalTok{current}\OperatorTok{,}\NormalTok{ points \} }\ImportTok{as}\NormalTok{ MedisContour)}\OperatorTok{;}
\NormalTok{      \}}
      \CommentTok{// Reset both points AND current so the next contour does not inherit}
      \CommentTok{// stale group/sliceDistance values from the previous block.}
\NormalTok{      current }\OperatorTok{=}\NormalTok{ \{\}}\OperatorTok{;}
\NormalTok{      points }\OperatorTok{=}\NormalTok{ []}\OperatorTok{;}
\NormalTok{    \} }\ControlFlowTok{else} \ControlFlowTok{if}\NormalTok{ (line}\OperatorTok{.}\FunctionTok{trim}\NormalTok{() }\OperatorTok{\&\&} \OperatorTok{!}\NormalTok{line}\OperatorTok{.}\FunctionTok{startsWith}\NormalTok{(}\StringTok{"\#"}\NormalTok{)) \{}
      \KeywordTok{const}\NormalTok{ [x}\OperatorTok{,}\NormalTok{ y}\OperatorTok{,}\NormalTok{ z] }\OperatorTok{=}\NormalTok{ line}\OperatorTok{.}\FunctionTok{trim}\NormalTok{()}\OperatorTok{.}\FunctionTok{split}\NormalTok{(}\SpecialStringTok{/}\SpecialCharTok{\textbackslash{}s+}\SpecialStringTok{/}\NormalTok{)}\OperatorTok{.}\FunctionTok{map}\NormalTok{(}\BuiltInTok{Number}\NormalTok{)}\OperatorTok{;}
      \ControlFlowTok{if}\NormalTok{ (}\OperatorTok{!}\PreprocessorTok{isNaN}\NormalTok{(x)) points}\OperatorTok{.}\FunctionTok{push}\NormalTok{(\{ x}\OperatorTok{,}\NormalTok{ y}\OperatorTok{,}\NormalTok{ z \})}\OperatorTok{;}
\NormalTok{    \}}
\NormalTok{  \}}
  \ControlFlowTok{if}\NormalTok{ (current}\OperatorTok{.}\AttributeTok{group} \OperatorTok{\&\&}\NormalTok{ points}\OperatorTok{.}\AttributeTok{length} \OperatorTok{\textgreater{}} \DecValTok{0}\NormalTok{) \{}
\NormalTok{    contours}\OperatorTok{.}\FunctionTok{push}\NormalTok{(\{ }\OperatorTok{...}\NormalTok{current}\OperatorTok{,}\NormalTok{ points \} }\ImportTok{as}\NormalTok{ MedisContour)}\OperatorTok{;}
\NormalTok{  \}}
  \ControlFlowTok{return}\NormalTok{ contours}\OperatorTok{;}
\NormalTok{\}}
\end{Highlighting}
\end{Shaded}

\subsubsection{Direct Client-Side Mesh Construction (\textless{} 100
ms)}\label{direct-client-side-mesh-construction-100-ms}

Contour rings are connected into a tube mesh directly in the browser ---
no backend round-trip needed. Algorithm: connect ring N to ring N+1 with
triangle pairs.

\textbf{Performance comparison:} \textbar{} Method \textbar{} Latency
\textbar{} Network \textbar{}
\textbar--------\textbar---------\textbar---------\textbar{} \textbar{}
Backend (buildstl.py → STL → download) \textbar{} 500--2000 ms
\textbar{} Required \textbar{} \textbar{} \textbf{Client-side (TXT →
NVMesh)} \textbar{} \textbf{50--100 ms} \textbar{} \textbf{None}
\textbar{}

\begin{center}\rule{0.5\linewidth}{0.5pt}\end{center}

\subsection{Mesh Generation
Strategies}\label{mesh-generation-strategies}

\subsubsection{Three Approaches}\label{three-approaches}

\begin{longtable}[]{@{}
  >{\raggedright\arraybackslash}p{(\columnwidth - 8\tabcolsep) * \real{0.2273}}
  >{\raggedright\arraybackslash}p{(\columnwidth - 8\tabcolsep) * \real{0.1818}}
  >{\raggedright\arraybackslash}p{(\columnwidth - 8\tabcolsep) * \real{0.1591}}
  >{\raggedright\arraybackslash}p{(\columnwidth - 8\tabcolsep) * \real{0.2045}}
  >{\raggedright\arraybackslash}p{(\columnwidth - 8\tabcolsep) * \real{0.2273}}@{}}
\toprule\noalign{}
\begin{minipage}[b]{\linewidth}\raggedright
Approach
\end{minipage} & \begin{minipage}[b]{\linewidth}\raggedright
Method
\end{minipage} & \begin{minipage}[b]{\linewidth}\raggedright
Speed
\end{minipage} & \begin{minipage}[b]{\linewidth}\raggedright
Quality
\end{minipage} & \begin{minipage}[b]{\linewidth}\raggedright
Use Case
\end{minipage} \\
\midrule\noalign{}
\endhead
\bottomrule\noalign{}
\endlastfoot
\textbf{Ultra-Simple} & Connect contour rings → STL & \textless{} 50 ms
& Faceted & MEDIS export \\
\textbf{Client-side vtk.js} & Marching cubes on mask & \textless{} 1 s &
Good & Interactive preview \\
\textbf{Server nii2mesh → MZ3} & Decimation + smoothing & 1--3 s & High
& Final visualisation \\
\end{longtable}

\subsubsection{Recommended Web Format:
MZ3}\label{recommended-web-format-mz3}

\begin{itemize}
\tightlist
\item
  3--5× smaller than PLY, 10× smaller than STL
\item
  Binary, gzip-compressed, native Niivue support
\item
  Target: \textless{} 5 MB per mesh, 50K--200K triangles, 60 FPS
  rendering
\end{itemize}

\begin{center}\rule{0.5\linewidth}{0.5pt}\end{center}

\subsection{Centerline Extraction \& Straightened MPR
(CPR)}\label{centerline-extraction-straightened-mpr-cpr}

\subsubsection{Overview}\label{overview}

\textbf{Straightened MPR} (Curved Planar Reformation) ``unfolds'' a
tortuous vessel into a straight view. Essential for assessing stenosis
and plaque distribution along the entire vessel length.

\textbf{Three steps:} 1. \textbf{Centerline extraction} --- centroid of
lumen contours (from MEDIS) or Voronoi skeletonisation (from AI mask) 2.
\textbf{Bishop (parallel transport) frame} --- compute Tangent (T) then
propagate Normal (N) and Binormal (B) without relying on curvature (see
§12.2 for why Frenet-Serret N must not be used) 3. \textbf{Cross-section
sampling} --- extract perpendicular slices → stack into straightened
volume

\subsubsection{Mathematical Foundation}\label{mathematical-foundation}

\textbf{Tangent} (finite differences):

\begin{verbatim}
T[i] = normalize(centerline[i+1] - centerline[i-1])  // central difference
\end{verbatim}

\textbf{Normal --- Bishop (parallel transport) frame:}

\begin{quote}
\textbf{Why not Frenet-Serret:} The Frenet-Serret principal normal
\texttt{N\ =\ normalize(dT/ds)} is undefined and flips 180° in
low-curvature (near-straight) vessel segments, producing rotating or
jumping cross-sections. Bishop frame (rotation-minimizing frame)
propagation\textsuperscript{6,7} avoids this by carrying an initial
normal forward without relying on curvature.
\end{quote}

\begin{verbatim}
// Initialise with any vector perpendicular to T[0]
N[0] = arbitrary_perpendicular(T[0])

// Propagate: project previous N onto the plane perpendicular to new T
for i in 1..n-1:
    N[i] = normalize(N[i-1] - dot(N[i-1], T[i]) * T[i])
\end{verbatim}

\textbf{Binormal}: \texttt{B{[}i{]}\ =\ T{[}i{]}\ ×\ N{[}i{]}}

\textbf{Cross-section point}:
\texttt{Q(u,v)\ =\ P{[}i{]}\ +\ u·N{[}i{]}\ +\ v·B{[}i{]}}

\subsubsection{Interactive Controls}\label{interactive-controls}

\begin{itemize}
\tightlist
\item
  \textbf{Position slider:} Select centerline point (0 → N-1)
\item
  \textbf{Rotation slider:} Rotate cross-section around T (Rodrigues'
  formula)
\item
  \textbf{Zoom slider:} Adjust cross-section FOV
\item
  \textbf{Quad-view:} CCTA overview \textbar{} cross-section \textbar{}
  straightened MPR \textbar{} 3D mesh
\end{itemize}

\begin{center}\rule{0.5\linewidth}{0.5pt}\end{center}

\subsection{Deployment, Performance \&
Security}\label{deployment-performance-security}

\subsubsection{Hardware Requirements}\label{hardware-requirements}

\paragraph{Development Environment}\label{development-environment}

\begin{longtable}[]{@{}
  >{\raggedright\arraybackslash}p{(\columnwidth - 6\tabcolsep) * \real{0.2292}}
  >{\raggedright\arraybackslash}p{(\columnwidth - 6\tabcolsep) * \real{0.1875}}
  >{\raggedright\arraybackslash}p{(\columnwidth - 6\tabcolsep) * \real{0.2708}}
  >{\raggedright\arraybackslash}p{(\columnwidth - 6\tabcolsep) * \real{0.3125}}@{}}
\toprule\noalign{}
\begin{minipage}[b]{\linewidth}\raggedright
Component
\end{minipage} & \begin{minipage}[b]{\linewidth}\raggedright
Minimum
\end{minipage} & \begin{minipage}[b]{\linewidth}\raggedright
Recommended
\end{minipage} & \begin{minipage}[b]{\linewidth}\raggedright
Current Setup
\end{minipage} \\
\midrule\noalign{}
\endhead
\bottomrule\noalign{}
\endlastfoot
\textbf{GPU} & RTX 3060 (8 GB) & RTX 2080 Ti (11 GB) & 2× RTX 2080 Ti
(11 GB each) \\
\textbf{CPU} & 6-core & 8+ core & Modern workstation \\
\textbf{RAM} & 16 GB & 32 GB+ & Sufficient \\
\textbf{Storage} & 500 GB SSD & 1 TB+ NVMe & Adequate \\
\textbf{CUDA} & 11.8+ & 12.2+ & CUDA 12.2 \\
\end{longtable}

\textbf{Development use cases:} model testing, single-user interactive
segmentation, DISCHARGE subset processing, active-learning experiments.

\paragraph{Production Environment (Charité Clinical
Deployment)}\label{production-environment-charituxe9-clinical-deployment}

\begin{longtable}[]{@{}
  >{\raggedright\arraybackslash}p{(\columnwidth - 6\tabcolsep) * \real{0.2750}}
  >{\raggedright\arraybackslash}p{(\columnwidth - 6\tabcolsep) * \real{0.2250}}
  >{\raggedright\arraybackslash}p{(\columnwidth - 6\tabcolsep) * \real{0.3250}}
  >{\raggedright\arraybackslash}p{(\columnwidth - 6\tabcolsep) * \real{0.1750}}@{}}
\toprule\noalign{}
\begin{minipage}[b]{\linewidth}\raggedright
Component
\end{minipage} & \begin{minipage}[b]{\linewidth}\raggedright
Minimum
\end{minipage} & \begin{minipage}[b]{\linewidth}\raggedright
Recommended
\end{minipage} & \begin{minipage}[b]{\linewidth}\raggedright
Notes
\end{minipage} \\
\midrule\noalign{}
\endhead
\bottomrule\noalign{}
\endlastfoot
\textbf{GPU Cluster} & 2× RTX 4090 (24 GB) & 4× A100/H100 (40--80 GB) &
Concurrent clinical users \\
\textbf{CPU} & 16-core & 32+ core & FastAPI + preprocessing \\
\textbf{RAM} & 64 GB & 128 GB+ & Multiple 3D volumes in memory \\
\textbf{Storage} & 10 TB+ & 50 TB+ & DISCHARGE + clinical data \\
\textbf{Network} & 10 Gbps & 25 Gbps+ & Fast volume transfers \\
\textbf{Redundancy} & RAID 10 & Distributed storage & Clinical data
safety \\
\end{longtable}

\textbf{Production requirements:} GDPR compliance (all data on-premise),
99.9 \% uptime, 5--10 concurrent radiologists, full DISCHARGE batch
processing, automated backup and failover.

\subsubsection{GDPR \& Deployment
Constraints}\label{gdpr-deployment-constraints}

\textbf{WARNING: IN-HOUSE BACKEND MANDATORY FOR CLINICAL USE}

For any clinical deployment, the SAM-Med3D-turbo model \textbf{MUST} run
on Charité's on-premise infrastructure:

\begin{longtable}[]{@{}
  >{\raggedright\arraybackslash}p{(\columnwidth - 4\tabcolsep) * \real{0.3636}}
  >{\raggedright\arraybackslash}p{(\columnwidth - 4\tabcolsep) * \real{0.2424}}
  >{\raggedright\arraybackslash}p{(\columnwidth - 4\tabcolsep) * \real{0.3939}}@{}}
\toprule\noalign{}
\begin{minipage}[b]{\linewidth}\raggedright
Requirement
\end{minipage} & \begin{minipage}[b]{\linewidth}\raggedright
Reason
\end{minipage} & \begin{minipage}[b]{\linewidth}\raggedright
Alternative
\end{minipage} \\
\midrule\noalign{}
\endhead
\bottomrule\noalign{}
\endlastfoot
\textbf{On-premise GPU servers} & GDPR --- patient data cannot leave
hospital & Not allowed: cloud inference \\
\textbf{Charité network} & Secure medical data transmission & Not
allowed: public internet \\
\textbf{Hospital authentication} & User access control and audit trails
& Not allowed: anonymous access \\
\textbf{Medical device certification} & Clinical safety and regulatory
compliance & Not allowed: research-only deployment \\
\end{longtable}

\begin{verbatim}
Browser (Hospital Network) → Charité Firewall → Internal GPU Cluster
       ↓                           ↓                        ↓
   UI/UX + Niivue            Load Balancer          SAM-Med3D-turbo
   3D visualization          + Authentication       PyTorch Inference
   User prompts              + Logging               + Medical Data Store
\end{verbatim}

\textbf{Non-clinical use cases (anonymised data allowed):} research
demos with cloud GPU; local development with sample datasets;
open-source contributions via GitHub with synthetic data; educational
browser-only UI with mock backend.

\subsubsection{Infrastructure}\label{infrastructure}

\begin{longtable}[]{@{}ll@{}}
\toprule\noalign{}
Component & Technology \\
\midrule\noalign{}
\endhead
\bottomrule\noalign{}
\endlastfoot
Containerisation & Docker + NVIDIA Container Toolkit \\
GPU & 4× NVIDIA A100 (80 GB each) \\
Embedding cache & Redis (sub-second for repeated prompts) \\
Batch queue & Celery + Redis broker \\
Authentication & LDAP / Charité SSO \\
Compliance & GDPR (all data on-premise, no cloud) \\
\end{longtable}

\subsubsection{Performance Targets}\label{performance-targets}

\begin{longtable}[]{@{}
  >{\raggedright\arraybackslash}p{(\columnwidth - 2\tabcolsep) * \real{0.5000}}
  >{\raggedright\arraybackslash}p{(\columnwidth - 2\tabcolsep) * \real{0.5000}}@{}}
\toprule\noalign{}
\begin{minipage}[b]{\linewidth}\raggedright
Metric
\end{minipage} & \begin{minipage}[b]{\linewidth}\raggedright
Target
\end{minipage} \\
\midrule\noalign{}
\endhead
\bottomrule\noalign{}
\endlastfoot
Single-click → mask & \textless{} 2 s end-to-end \\
Embedding computation (first prompt) & \textasciitilde1 s \\
Subsequent prompts (cached embedding) & \textless{} 0.5 s \\
Mesh generation (MZ3) & \textless{} 3 s \\
Batch throughput & \textasciitilde50 cases / hour (4 GPUs) ---
bottleneck is DICOM→NIfTI preprocessing, not GPU \\
CTA volume memory & \textasciitilde268 MB (512³ × 2 bytes = 268,435,456
bytes) \\
\end{longtable}

\subsubsection{Memory Budget}\label{memory-budget}

\begin{longtable}[]{@{}
  >{\raggedright\arraybackslash}p{(\columnwidth - 2\tabcolsep) * \real{0.6471}}
  >{\raggedright\arraybackslash}p{(\columnwidth - 2\tabcolsep) * \real{0.3529}}@{}}
\toprule\noalign{}
\begin{minipage}[b]{\linewidth}\raggedright
Component
\end{minipage} & \begin{minipage}[b]{\linewidth}\raggedright
Size
\end{minipage} \\
\midrule\noalign{}
\endhead
\bottomrule\noalign{}
\endlastfoot
CTA volume (512³, int16) & \textasciitilde268 MB \\
SAM-Med3D-turbo (FP16) & \textasciitilde4 GB VRAM \\
Embedding cache (per volume) & \textasciitilde50--500 MB (estimate ---
depends on patch count and feature map size; to be measured) \\
Straightened volume (64² × 200) & \textasciitilde8 MB \\
Mesh (MZ3, per vessel) & \textless{} 5 MB \\
\end{longtable}

\begin{center}\rule{0.5\linewidth}{0.5pt}\end{center}

\subsection{Research Roadmap \&
Milestones}\label{research-roadmap-milestones}

\subsubsection{Phase 1: MVP --- MEDIS TXT Visualisation (Weeks
1--4)}\label{phase-1-mvp-medis-txt-visualisation-weeks-14}

\begin{itemize}
\tightlist
\item[$\square$]
  MEDIS TXT parser + client-side mesh in Niivue
\item[$\square$]
  Centerline extraction + straightened MPR with \textbf{Bishop (parallel
  transport) frame} (not Frenet-Serret --- see §12.2)
\item[$\square$]
  Quad-view layout with interactive sliders
\item[$\square$]
  NIfTI.gz / DICOM loading
\end{itemize}

\subsubsection{Phase 2: SAM-Med3D Integration (Weeks
5--8)}\label{phase-2-sam-med3d-integration-weeks-58}

\begin{itemize}
\tightlist
\item[$\square$]
  Backend: load turbo checkpoint, expose \texttt{/api/segment/point}
\item[$\square$]
  Frontend: click → prompt → mask overlay
\item[$\square$]
  Redis embedding cache for sub-second repeated prompts
\item[$\square$]
  \textbf{Implement \texttt{dense\_prompt} support} --- either (a)
  custom prompt encoder modification to accept a prior mask, or (b)
  lumen mask surface → additional positive point prompts; required
  blocker for dual-wall pipeline
\item[$\square$]
  Dual-wall sequential segmentation pipeline (depends on above)
\end{itemize}

\subsubsection{Phase 3: DISCHARGE Evaluation (Weeks
9--12)}\label{phase-3-discharge-evaluation-weeks-912}

\begin{itemize}
\tightlist
\item[$\square$]
  \textbf{Zero-shot baseline --- myocardium first:} \texttt{LV\_MYO} is
  the closest match to ACDC pre-training targets; establish Dice
  baseline here before coronary lumen (see §2.5 strategic implication
  note)
\item[$\square$]
  Zero-shot baseline on coronary lumen (held-out DISCHARGE cases)
\item[$\square$]
  Quantify Dice, Hausdorff, diameter stenosis \% correlation vs.~MEDIS
  QAngio CT expert measurements
\item[$\square$]
  Identify failure modes (motion, calcification, bifurcations,
  patch-boundary seams)
\end{itemize}

\subsubsection{Phase 4: Fine-tuning \& Active Learning (Weeks
13--20)}\label{phase-4-fine-tuning-active-learning-weeks-1320}

\begin{itemize}
\tightlist
\item[$\square$]
  Fine-tune SAM-Med3D on DISCHARGE annotations (nnU-Net-style data prep)
\item[$\square$]
  Active-learning loop\textsuperscript{8}: expert corrections → weekly
  re-training. \textbf{Correction data model:} a ``correction'' is a
  full per-structure binary NIfTI mask saved by the radiologist after
  local brush edits in the browser. The backend must convert the
  in-browser voxel edits (sparse difference from AI mask) into a
  complete corrected NIfTI matching §2.9 training format before
  ingestion into the fine-tuning pipeline.
\item[$\square$]
  Benchmark against task-specific nnU-Net baseline
\end{itemize}

\subsubsection{Phase 5: Prostate Extension (Weeks
21--28)}\label{phase-5-prostate-extension-weeks-2128}

\begin{itemize}
\tightlist
\item[$\square$]
  Adapt pipeline for prostate mpMRI (multi-sequence input)
\item[$\square$]
  Zone-specific class labels (PZ / TZ / lesion)
\item[$\square$]
  Validate on Charité prostate cohort
\end{itemize}

\subsubsection{Phase 6: Clinical Validation \& Publication (Weeks
29--36)}\label{phase-6-clinical-validation-publication-weeks-2936}

\begin{itemize}
\tightlist
\item[$\square$]
  Prospective reader study (Dice vs.~time vs.~inter-observer)
\item[$\square$]
  Open-source release + MedSegFM competition baseline
\item[$\square$]
  Publication: ``Foundation-model-assisted coronary CTA segmentation at
  scale''
\end{itemize}

\subsubsection{Quarterly Research
Milestones}\label{quarterly-research-milestones}

\begin{longtable}[]{@{}
  >{\raggedright\arraybackslash}p{(\columnwidth - 4\tabcolsep) * \real{0.3214}}
  >{\raggedright\arraybackslash}p{(\columnwidth - 4\tabcolsep) * \real{0.3929}}
  >{\raggedright\arraybackslash}p{(\columnwidth - 4\tabcolsep) * \real{0.2857}}@{}}
\toprule\noalign{}
\begin{minipage}[b]{\linewidth}\raggedright
Quarter
\end{minipage} & \begin{minipage}[b]{\linewidth}\raggedright
Milestone
\end{minipage} & \begin{minipage}[b]{\linewidth}\raggedright
Status
\end{minipage} \\
\midrule\noalign{}
\endhead
\bottomrule\noalign{}
\endlastfoot
Q1 2026 (ends 2026-03-31) & Zero-shot baseline on DISCHARGE + MVP
deployed & In progress - update before end of quarter \\
Q2 2026 & Active-learning loop running; Dice ≥ 0.80 on coronary lumen &
Pending \\
Q3 2026 & Prospective reader study; prostate pipeline validated &
Pending \\
Q4 2026 & Open-source release; conference/journal submission &
Pending \\
\end{longtable}

\begin{center}\rule{0.5\linewidth}{0.5pt}\end{center}

\subsection{Related Medical Segmentation
Models}\label{related-medical-segmentation-models}

\begin{longtable}[]{@{}
  >{\raggedright\arraybackslash}p{(\columnwidth - 8\tabcolsep) * \real{0.1029}}
  >{\raggedright\arraybackslash}p{(\columnwidth - 8\tabcolsep) * \real{0.1618}}
  >{\raggedright\arraybackslash}p{(\columnwidth - 8\tabcolsep) * \real{0.1618}}
  >{\raggedright\arraybackslash}p{(\columnwidth - 8\tabcolsep) * \real{0.1324}}
  >{\raggedright\arraybackslash}p{(\columnwidth - 8\tabcolsep) * \real{0.4412}}@{}}
\toprule\noalign{}
\begin{minipage}[b]{\linewidth}\raggedright
Model
\end{minipage} & \begin{minipage}[b]{\linewidth}\raggedright
Dimension
\end{minipage} & \begin{minipage}[b]{\linewidth}\raggedright
Modalities
\end{minipage} & \begin{minipage}[b]{\linewidth}\raggedright
Prompts
\end{minipage} & \begin{minipage}[b]{\linewidth}\raggedright
Key Difference from SAM-Med3D
\end{minipage} \\
\midrule\noalign{}
\endhead
\bottomrule\noalign{}
\endlastfoot
SAM (Meta) & 2D & Natural images & Point/box/text & No medical training,
2D only \\
SAM-Med2D & 2D & Medical (slice-wise) & Point/box & 2D → cannot capture
volumetric context \\
MedSAM & 2D & Medical (slice-wise) & Box only & Simpler architecture,
box prompts only \\
SAM-Med3D & \textbf{3D} & \textbf{Medical (volumetric)} & \textbf{3D
point} & \textbf{Native 3D --- our choice} \\
nnU-Net & 3D & Medical (task-specific) & None (automatic) & Not
promptable; requires per-task training \\
\end{longtable}

\begin{center}\rule{0.5\linewidth}{0.5pt}\end{center}

\emph{This document is the living foundation of the Charité Segment
Platform research project. All claims about SAM-Med3D are verified
against the published paper (arXiv:2310.15161), official GitHub README,
and HuggingFace model/dataset cards. Critical caveats about zero-shot
performance on coronary CTA and prostate mpMRI (neither directly
evaluated in the paper) are noted throughout. Full bibliography in
\texttt{plan/references.bib}.}

\section*{References}\label{bibliography}
\addcontentsline{toc}{section}{References}

\phantomsection\label{refs}
\begin{CSLReferences}{0}{1}
\bibitem[\citeproctext]{ref-Zhang2024SAMMed3D}
\CSLLeftMargin{1. }%
\CSLRightInline{Zhang H, Xia Y, Zhao S, Luo M, Zhang Q, Xu M, et al.
SAM-Med3D: Towards general-purpose segmentation models for volumetric
medical images. In: European conference on computer vision (ECCV)
workshops: Bioimage computing {[}Internet{]}. Springer; 2024. p. 1--18.
doi:\href{https://doi.org/10.1007/978-3-031-72469-6_1}{10.1007/978-3-031-72469-6\_1}}

\bibitem[\citeproctext]{ref-Isensee2021nnUNet}
\CSLLeftMargin{2. }%
\CSLRightInline{Isensee F, Jaeger PF, Kohl SAA, Petersen J, Maier-Hein
KH. nnU-net: A self-configuring method for deep learning-based
biomedical image segmentation. Nature Methods. 2021;18(2):203--11.
doi:\href{https://doi.org/10.1038/s41592-020-01008-z}{10.1038/s41592-020-01008-z}}

\bibitem[\citeproctext]{ref-DISCHARGE2022NEJM}
\CSLLeftMargin{3. }%
\CSLRightInline{DISCHARGE Trial Group. Computed tomography or invasive
coronary angiography in patients with stable chest pain. New England
Journal of Medicine. 2022;386(19):1795--806.
doi:\href{https://doi.org/10.1056/NEJMoa2200963}{10.1056/NEJMoa2200963}}

\bibitem[\citeproctext]{ref-Maurovich2020CADRADS}
\CSLLeftMargin{4. }%
\CSLRightInline{Maurovich-Horvat P, Ferencik M, Voros S, Merkely B,
Hoffmann U. Comprehensive plaque assessment by coronary CT angiography.
Nature Reviews Cardiology. 2014;11(7):390--402.
doi:\href{https://doi.org/10.1038/nrcardio.2014.60}{10.1038/nrcardio.2014.60}}

\bibitem[\citeproctext]{ref-Niivue2024}
\CSLLeftMargin{5. }%
\CSLRightInline{Rorden C, Taylor H, Gallivan K, Hanayik T. Niivue: A
WebGL2-based medical image viewer {[}Internet{]}. GitHub; 2024.
doi:\href{https://doi.org/10.5281/zenodo.8320937}{10.5281/zenodo.8320937}}

\bibitem[\citeproctext]{ref-Wang2008RotationMinimizing}
\CSLLeftMargin{6. }%
\CSLRightInline{Wang W, Jüttler B, Zheng Y, Liu Y. Computation of
rotation minimizing frames. ACM Transactions on Graphics.
2008;27(1):1--18.
doi:\href{https://doi.org/10.1145/1330511.1330512}{10.1145/1330511.1330512}}

\bibitem[\citeproctext]{ref-Kanitsar2002CPR}
\CSLLeftMargin{7. }%
\CSLRightInline{Kanitsar A, Ropinski T, Fleischmann D, König A, Gröller
ME. CPR - curved planar reformation. In: IEEE visualization. IEEE; 2002.
p. 413--20.
doi:\href{https://doi.org/10.1109/VIS.2002.5930066}{10.1109/VIS.2002.5930066}}

\bibitem[\citeproctext]{ref-Budd2021ActiveLearning}
\CSLLeftMargin{8. }%
\CSLRightInline{Budd S, Robinson EC, Kainz B. A survey on active
learning and human-in-the-loop deep learning for medical image analysis.
Medical Image Analysis. 2021;71:102062.
doi:\href{https://doi.org/10.1016/j.media.2021.102062}{10.1016/j.media.2021.102062}}

\end{CSLReferences}

\end{document}
